\documentclass[../main.tex]{subfiles}
\begin{document}
%==========================================TIÊU ĐỀ TẬP========================================
\begin{table}[h]
  \begin{adjustwidth}{-1cm}{}
    \begin{tabular}{>{\centering\arraybackslash}p{6cm}|>{\raggedright\arraybackslash}p{12.5cm}}
      \multirow{5}{6cm}
      % Cột bên trái: ÉPISODE và số tập
      {\\[-40pt]
        \flaregothic\fontsize{20pt}{20pt}\selectfont \centering
        {\contour{errorthree}{\textcolor{black}{ÉPISODE}}}\color{black}\\
        \burbank\fontsize{72pt}{72pt}\selectfont
        \includegraphics[height=3.12359cm]{../../../../Image/Book?/number/num16.png}
        \\[18pt]
        % \flaregothic\fontsize{14pt}{14pt}\selectfont
        % \color{epattr}BIENVENUE, \uline{\color{Banpire} Mel} !
        % \flaregothic\fontsize{18pt}{18pt}\selectfont
        % \color{epattr}ÉPISODE PERDU
        \color{black}
      }
       &
      % Dòng thứ nhất: tên tập tiếng Nhật
      {\fotskip\fontsize{18pt}{18pt}\selectfont \ruby{本}{ほん}\ruby{当}{とう}の\ruby{私}{わたし}
      } \\[6pt]
      % Dòng thứ hai: tên tập tiếng Anh
       & {\cafeteria\fontsize{18pt}{18pt}\selectfont Who I Really Am}                 \\[4pt]
      % Dòng thứ ba: tên tập tiếng Việt
       & {\vnmsans\fontsize{18pt}{18pt}\selectfont Tôi\dots\ là ai?}                  \\[8pt]
      % Dòng thứ tư: Nhân vật xuất hiện
       & {\caviar{\fontsize{14pt}{14pt}\selectfont Track used: The War}}
      \\[6pt]
      % Dòng cuối cùng: Thời gian diễn ra của tập (thường sẽ là ngày phát hành)
       & {\continuum\fontsize{16pt}{16pt}\selectfont Ngày phát hành: 04/06/2022}      \\[8pt]
    \end{tabular}
  \end{adjustwidth}
\end{table}
%===========================================NỘI DUNG==========================================
\color{black}

\begin{center}
  ``Đã đến lúc cậu phải đối diện với sự thật rồi, Sora.''
\end{center}

Lời nói ấy vang vọng như một nhát búa nện thẳng vào lồng ngực. Sora đứng sững, đôi môi run rẩy, ánh mắt rỗng hoác. Cô lắp bắp, giọng đứt đoạn, gần như nghẹn lại: ``Đ-Đối diện\dots\ với cái gì cơ\dots?''

Kyuen nheo mắt lại, đồng tử co rút; cậu tiến nửa bước về phía Sora rồi khựng lại, bàn tay giật giật như muốn chạm vào vai cô để níu kéo nhưng lại sợ vỡ mất thứ gì đang lung lay.

Saikyou siết chặt tay đến trắng xương, khớp ngón kêu răng rắc; cô nhớ lại lời Shino căn dặn trước khi đặt chân vào vòng lặp này: ``Khi vòng lặp lộ diện, đừng để nó nuốt chửng cả bạn lẫn mình.''

Reiko ôm lấy cánh tay mình, khuôn mặt cô căng ra, chờ xem những gì diễn ra tiếp theo.

Shino mở mắt to, cái nhìn xoáy thẳng vào Sora. Không phải ngạc nhiên, mà là một sự khẳng định lạnh lùng. Cô nói chậm rãi, từng chữ như rơi xuống nặng nề: ``Cậu vẫn không chịu giác ngộ à? Hãy nhìn những người xung quanh cậu đi. Thế giới lý tưởng hão huyền đó\dots\ chính cậu đã tạo ra nó đấy.''

\img{e2-8a}

Sora chớp mắt, quay phắt ra phía sau. Cơn lạnh buốt chạy dọc sống lưng, cô thấy mình như đang lạc vào một bức tranh bị rút hết màu: cô vừa nhìn thấy Akane và Uzuki đang cúi gập xuống, miệng lẩm bẩm nói những lời vô nghĩa, khuôn mặt dần mờ nhạt như bóng đen đang bò ra từ trong lòng họ.

Trông thấy cảnh tượng trước mắt, Kyuen thoáng giật mình, vội lùi một bước, lắp bắp: ``Có vẻ như\dots\ họ đang biến thành `chúng' rồi\dots''

\img{e2-8b}

Cô nàng tai chó đứng gục đầu, đôi mắt mất dần ánh sáng, lờ đờ như cá ươn. Cô nàng tai thỏ đứng cạnh, mắt trống rỗng, môi hé mở lặp đi lặp lại những lời xa xưa như tiếng vọng méo mó.

``Phải quay về thôi\dots\ Đi đứng cẩn thận vào chứ. Cứ chạy như thế kẻo lại làm thương tích ai đó đấy,'' giọng cô nàng bờm thỏ vang lên, chính xác là câu cô từng thốt ra khi Reiko vô tình va vào mình khi đang chuẩn bị lên tàu.

``Phải quay về đó thôi\dots\ Ừ nhỉ. Mình cũng thấy tò mò đấy. Bọn mình cũng muốn các cậu dẫn dắt tham quan vòng quanh thị trấn này,'' cô nàng tai chó thì thào, ánh mắt vô hồn như bị hút về khoảng trống vô tận.

Từ phía sau, Yuka khẽ cất tiếng cười lanh lảnh, nhưng không hề còn là con người bình thường: ``Ta cùng quay về đó thôi nhé\~{} Quay về Aogami. Sắp tới chúng ta sẽ là hàng xóm của nhau rồi hen\~{}''

\begin{figure}[H]
  \centering
  \begin{subfigure}[t]{0.45\textwidth}
    \centering
    \includegraphics[width=8cm]{../../../../Image/Book?/e2-8c1.png}
  \end{subfigure}
  \quad
  \begin{subfigure}[t]{0.45\textwidth}
    \centering
    \includegraphics[width=8cm]{../../../../Image/Book?/e2-8c2.png}
  \end{subfigure}
\end{figure}

Kyuen đứng chết lặng, đầu gối run nhè nhẹ, cảm thấy mình như đang chìm vào bùn.  ``Không thể nhầm lẫn vào đâu được\dots'' cậu thều thào.

Saikyou rùng mình, bật ra một tiếng rít qua kẽ răng: ``Chết tiệt\dots\ đúng là nó rồi.''

Reiko cắn môi tới bật máu, đôi mắt chao đảo, không thể tin nổi cả ba người bạn ấy cũng đang trượt dần vào bóng tối vô hồn: ``Không lẽ\dots\ ba người đó cũng sẽ có số phận như vậy?''

Sora thụt lùi lại, hai tay ôm lấy đầu, hơi thở gấp gáp. Mồ hôi vã ra như tắm, trượt dọc gò má. Cô thấy những mảnh ký ức đang vỡ ra từng mảnh như gương rơi: nụ cười Akane khi trao cô chiếc bánh tai yêu, cái gật đầu thân thiện của Uzuki, tiếng hát Yuka vang trong gió chiều.

``D-Dừng lại đi mà\dots\ Akane-san\dots\ Uzuki-san\dots\ Yuka-san\dots\ tại sao lại như vậy?'' giọng cô vỡ vụn. ``Tôi\dots\ Tôi chỉ muốn làm bạn với mọi người thôi mà\dots''

\img{e2-8d}

Những lời ấy bật ra như tiếng kêu cứu cuối cùng trong tuyệt vọng. Nhưng đối diện với cô, ba khuôn mặt bạn bè chỉ còn lại cái vỏ rỗng, lặp đi lặp lại những lời vô nghĩa. Cả thế giới mà Sora cố công xây đắp đang rạn nứt trước mắt, sụp đổ từng mảnh, và tất cả chỉ còn lại nỗi dằn vặt nhấn chìm cô trong nỗi sợ hãi không lối thoát.

Sora vẫn ngồi đó, hai tay run rẩy ôm chặt lấy đầu, hơi thở dồn dập. Còn Shino thì đứng thẳng lưng, ánh mắt như xuyên thấu qua lớp vỏ mà bản ngã mắt nâu đang cố gắng che chắn.

``Cậu vẫn định sẽ chạy trốn sao?'' cô cất giọng, từng chữ nặng nề rơi xuống, lạnh tanh và trầm chậm như dao cứa. ``Lặp lại tất cả mọi thứ, tiếp tục bám víu vào cái ảo tưởng đó?''

``Ể\dots?'' Sora ngẩng lên, đôi mắt dại đi vì kinh ngạc và sợ hãi.

Từ phía xa, Reiko khẽ lắc đầu. Giọng cô vang lên, không to, nhưng rành rọt đủ để cắt qua bầu không khí ngột ngạt: ``Mình nghĩ\dots\ cậu nên tỉnh mộng đi, Sora-san. Đừng cố nắm lấy một giấc mơ đã nát vụn nữa. Đã đến lúc cậu phải chấp nhận sự thật rồi.''

Sora nhìn cô, môi run run, rồi lại quay sang Shino. Nhưng thay vì buông bỏ, ánh mắt cô vẫn còn ánh lên tia níu kéo tuyệt vọng, như muốn bám víu lần cuối vào mảnh vụn lý tưởng đã rạn nứt.

Shino cau mày. Đôi mắt xanh vốn sắc lạnh giờ hằn rõ sự bất lực. ``Cậu vẫn ngang bướng như mọi khi nhỉ.''

Lời ấy nặng trĩu, đủ khiến Sora sững người, trái tim cô chao đảo.

Không chần chừ thêm, cô bước thẳng tới cuối toa tàu. Cánh cửa kim loại bật mở, tiếng gió gào rít ùa vào cùng tiếng xình xịch dồn dập của bánh sắt. Con tàu vẫn lao đi vun vút, rung lắc dữ dội.

\img{e2-8e}

Cô ngoái đầu lại, ánh mắt khắc nghiệt, buông thêm một nhát chém chí mạng: ``Cậu vẫn sẽ mãi chạy trốn thôi. Cứ chạy trốn, chạy trốn\dots\ và mãi mãi vẫn vậy.''

Sora cứng đờ, đôi môi mấp máy như người sắp ngạt thở. ``Tôi\dots\ Tôi đang chạy trốn ư?'' giọng cô nghẹn lại, như thể chính từ ấy chưa từng tồn tại trong từ điển của mình.

Shino không đáp. Cô chỉ xoay người, tiếp tục bước ra cánh cửa, bóng lưng cứng cỏi đối diện với cơn gió lùa.

``Đ-Đợi đã! Ý cậu là sao chứ!?'' Sora bật thốt, giọng lạc đi vì tuyệt vọng.

Ở phía sau, Kyuen khẽ nhắm mắt lại, buông một tiếng thở dài, như đã đoán trước tình huống này từ lâu. Saikyou chỉ lặng im quan sát, ánh mắt đanh lại, trong khi Reiko vẫn không rời ánh nhìn khỏi Sora, tựa hồ chờ xem cô bạn kia có dám đối diện thật sự hay không.

Shino khựng lại khi nghe tiếng Sora. Cô quay đầu, ánh mắt lạnh lùng quét qua, như thể chờ đợi điều này từ lâu. Sora bước lên một bước, định lao về phía Shino, nhưng ngay lập tức bị Uzuki, Akane và Yuka giữ chặt hai cánh tay, cả cơ thể bị níu lại.

``Quay về đi\dots'' Uzuki thì thầm.

``Hãy quay về đi\dots'' Akane và Yuka lặp lại, giọng mỏng manh, yếu ớt như vọng ra từ xa xăm.

\img{e2-8f}

Saikyou chứng kiến cảnh tượng đó, bàn tay bất giác run lên. Hình ảnh Kaoru níu lấy ngăn ba anh chị em cô tới ga Aogami bùng lên trong đầu cô: ánh mắt Kaoru, lạnh lẽo và cứng rắn như viên đạn xuyên qua ký ức, ép buộc cô phải đứng yên, bất lực. Giờ đây, cô thấy lại ánh mắt đó, nhưng lần này không phải nhắm vào mình, mà là vào Sora: một Sora đang tan rã đến thê thảm.

``Tại sao\dots\ Tại sao mọi người lại như vậy\dots?'' Sora hoảng loạn, giãy giụa, giọng lạc đi.

Kyuen, Saikyou và Reiko vẫn im lặng, chỉ nhìn. Bầu không khí trong toa tàu như bị nén chặt, mỗi tiếng động nhỏ cũng bị nuốt chửng. Họ hướng ánh mắt về phía Shino, như chờ một lời giải đáp. Shino chỉ đáp lại bằng một ánh nhìn, ngắn gọn mà sắc lạnh: ``Cứ để cô ấy giằng co như vậy đi.''

Trong đầu Sora, tất cả vỡ vụn thành một mớ hỗn loạn.

``\textit{Cái gì mới là thật?}'' cô tự hỏi, mồ hôi lạnh trượt dài. ``\textit{Mình đã tin vào điều gì? Những khuôn mặt thân thuộc này\dots\ chỉ là ảo ảnh ư? Vậy thì cái thế giới mà mình cố giữ lại, nơi mọi người vẫn còn sống, nơi mình không phải mất đi ai hết\dots\ cũng chỉ là giả dối?}''

Cảm giác ngột ngạt siết lấy ngực cô. Một phần trong cô muốn hét lên, muốn gào hỏi Shino cho rõ ràng: \textit{Tại sao phải phá vỡ nó? Tại sao không thể để mình tiếp tục sống trong giấc mơ ấy? Nhưng một phần khác lại rỉ máu, rít lên từng chữ: Nếu tất cả chỉ là mơ\dots\ thì bám víu để làm gì?}

Đôi mắt cô đỏ hoe, lập lòe trong nước mắt: ``\textit{Mình\dots\ không muốn\dots\ Mình không muốn tiếp tục mất đi bất cứ ai nữa. Nhưng nếu đây chỉ là ảo mộng\dots\ thì mình đã lừa dối chính bản thân. Mình sẽ mãi mãi chạy trốn\dots\ đúng như Shino nói sao?}''

Tiếng thì thầm ``quay về đi'' lặp lại, dồn dập như hàng trăm chiếc kim chích vào tâm trí. Sora hét lên một tiếng, rồi bất ngờ giật mạnh, hất văng Uzuki, Akane và Yuka. Ba thân thể kia ngã xuống, biến dạng, vỡ vụn như những con rối rách nát, không còn chút hơi thở nào.

``\textbf{Bỏ tôi ra!!!}''

\img{e2-8g}

Sora khụy gối, thở dốc, nhưng trong mắt đã không còn sự mơ màng nữa.

Shino khẽ mỉm cười, nụ cười hiếm hoi nhưng lạ lùng mãn nguyện. ``Cuối cùng\dots\ cậu cũng chịu rồi,'' cô thì thầm, như thốt ra cho chính mình nghe.

Ba anh chị em Haragen quan sát. Người anh trai siết chặt nắm tay, lòng dâng lên một cảm giác kỳ lạ, vừa nhẹ nhõm vì Sora đã thoát, vừa lạnh sống lưng vì những gì đang tới. Cô em út lại rùng mình, ý nghĩ về định đanh chồng chéo của chính họ bất giác ùa về. Người chị thì nuốt khan, chẳng dám rời mắt khỏi Shino.

``Mình sẽ giải quyết chuyện này từ đây,'' Shino lên tiếng, giọng cô vang đều trong toa. ``Mình sẽ đưa bản ngã `Sora' tới gần hơn với sự thật. Cho cô ta biết\dots\ cô ta thực sự là ai.''

*OE\dots\ OE\dots\ OE\dots* *XÌNH XỊCH, XÌNH XỊCH* *LỤC CỤC, LỤC CỤC*

Ngay khi câu nói vừa dứt, tiếng còi báo động chát chúa vang lên, đỏ rực đèn nhấp nháy khắp toa. Con tàu rùng mạnh, tăng tốc đến mức ghê rợn.

Shino liếc sang ba anh chị em, nói thêm: ``Các cậu cũng sẽ sớm được biết thân phận của mình thôi. Đừng lo\dots\ chúng ta sẽ còn gặp lại nhau nữa đấy.''

Sora ngẩng đầu, nhìn về ba thân thể đã hóa rối, há hốc mồm: ``Mọi người\dots\ không\dots\ không thể nào\dots'' Nhưng rồi cô bước đến Shino, đôi mắt đầy tuyệt vọng.

``Cậu định nói với tôi điều gì? Giải thích cho tôi đi! Tôi đã làm gì sai!?''

\img{e2-8i}

Shino đặt tay lên cửa tàu, từng bước đi ra ngoài, chỉ đáp lại ngắn gọn: ``Cứ đi theo tôi, rồi cậu sẽ hiểu.''

Cô nhảy xuống, bóng dáng tan biến vào làn sương đỏ. Trước khi biến mất, Shino ngoái lại, khẽ cười hiền: ``Bảo trọng nhé\dots

\begin{center}
  Hikawa Kuon-san, Hikawa Kaigou-san, và\dots\ Hikawa Suzuki-san.''
\end{center}

Ba anh em chết lặng. Tiếng gọi tên ấy như mũi dao cắm sâu vào tim họ.

``Những cái tên ấy\dots'' Kyuen thì thầm, giọng khô khốc. Trong đầu cậu, những mảnh ký ức mơ hồ về một căn nhà nhỏ ở Aogami, về tiếng gọi ``Hikawa'' vang lên trong gió, bất chợt ùa về. Cậu không biết mình là ai nữa, là Kyuen, hay là Kuon? Hay cả hai đều chỉ là tên của một con ma đã chết từ lâu?

``Không\dots\ Không thể nào\dots'' Saikyou lắc đầu, nhưng mặt trắng bệch. Đôi mắt cô trống rỗng, như thể vừa nhìn thấy một bức màn bị xé toạc. ``Kaigou\dots'' Cô lẩm bẩm tên ấy, và trong phút chốc, cô thấy mình đứng giữa một ga tàu khác, trong một thân xác khác, với một cái tên khác. Cô đã từng nghĩ mình là Saikyou, nhưng nếu cái tên ấy chỉ là lớp vỏ bọc cho một ký ức bị xóa sạch thì sao?

``Nếu vậy\dots\ thì chúng ta\dots\ thực sự là họ sao? Hay chỉ là\dots'' Reiko bỏ lửng, chìm trong tiếng còi hú. Cô không dám nghĩ tiếp. Cô sợ rằng nếu mình là Suzuki, thì cô bé Reiko với mái tóc cột hai bên và nụ cười nhút nhát kia chỉ là một vai diễn trong vở kịch của kẻ nào đó. Cô sợ rằng cả ba người họ chưa từng tồn tại. Chỉ là những bóng ma đang cố nhớ lại tên mình.

*OE\dots\ OE\dots\ OE\dots*

Đèn đỏ nhấp nháy, tiếng báo động dồn dập, đưa họ về thực tại: họ đang ở trên con tàu lao đi với vận tốc kinh hoàng. Trong toa, chỉ còn ba người họ, những kẻ duy nhất vẫn còn\dots\ sống.

Hay chỉ là\dots\ những kẻ chưa chịu chết?

\ct

Tiếng còi tàu rền vang, kéo dài như một lời rên rỉ xé toạc không gian, hòa lẫn cùng tiếng sắt thép nghiến ken két dưới gầm bánh xe. Mỗi nhịp xình xịch lại dội vào vách toa như những nhát búa đập thẳng vào màng tai, khiến người nghe khó phân biệt đâu là âm thanh của máy móc, đâu là tiếng than khóc của một linh hồn bị mắc kẹt.

Đèn trần lập lòe, ánh sáng trắng tái nhợt nhấp nháy từng hồi, chiếu lên những gương mặt xanh xao và đôi mắt căng thẳng, để rồi ngay sau đó lại nhấn chìm tất cả vào thứ bóng tối mơ hồ. Thỉnh thoảng, còi báo động rú lên một tràng, chói gắt đến mức cả không gian như rung bần bật, buộc những hành khách vô thức co rúm người lại, lòng bàn tay rịn mồ hôi lạnh.

Không ai nói gì, nhưng sự im lặng ấy không hề nhẹ nhõm, trái lại, nó lại ngột ngạt, tới mức không ai thở nổi.

Saikyou lao về phía cánh cửa cuối toa, nơi mà Shino và Sora đã nhảy ra ban nãy. Nhưng khi vừa với tay đến, cô nhận ra cửa đã bị khóa chặt từ bao giờ, bề mặt thép lạnh băng không hé mở lấy một khe hở.

``Chết tiệt\dots!'' cô đấm mạnh vào cánh cửa, tiếng kim loại vang lên chát chúa. ``Bị chặn rồi\dots! Không thể đuổi theo họ được nữa!''

Kyuen nghiến răng, mắt lia khắp toa tàu. ``Vậy thì phải tìm đường khác thôi\dots\ Có lẽ đi qua toa bên cạnh, hoặc tìm lối thoát hiểm!''

Reiko luống cuống: ``Nhưng\dots\ nếu chúng ta rời khỏi đây\dots\ thì họ\dots''

Cô ngoái đầu lại. Cùng lúc đó, hai anh chị lớn cũng quay lại, ba ánh mắt đồng loạt dừng lại ở phía cuối toa. Yuka, Akane và Uzuki vẫn đứng đó, bất động như tượng sáp, mắt trống rỗng nhìn về phía trước. Không ai nói gì, không ai cử động. Không khí như ngưng đọng.

\img{e2-8h}

``Họ\dots\ đang làm gì vậy\dots?'' Cô em út lẩm bẩm, giọng run lên từng chữ.

Kyuen bước lùi một bước, mắt không rời khỏi ba bóng dáng kia. ``Tránh xa họ ra, Reiko. Họ không còn là con người nữa.''

Chỉ trong chớp mắt, thân hình Yuka khẽ run lên, một cái giật nhẹ, như thể có dòng điện chạy qua. Rồi Akane cũng vậy. Uzuki cũng thế. Ba cơ thể đồng loạt co giật, nhịp độ càng lúc càng nhanh, càng lúc càng dữ dội.

``C-Cái quái gì thế này!?'' Saikyou bật thốt, lùi về phía sau.

Ba bóng người rung lên như những bức ảnh nhiễu sóng. Đường nét khuôn mặt bắt đầu méo mó, tan loãng thành những mảng đen sì, như mực loang trong nước. Không có máu, không có tiếng kêu, chỉ có sự biến mất từ từ, im lặng đến rợn người.

``Họ\dots\ đang bị xóa\dots'' Reiko nghẹn giọng, hai tay bịt miệng, mắt mở to đến trợn trừng.

Kyuen siết chặt tay, môi mấp máy: ``Không chỉ là biến mất\dots\ họ đang bị vòng lặp \emph{thu hồi}\dots\ như thể\dots\ như thể họ chỉ là những con rối đã hết tác dụng\dots''

Chỉ vài giây sau, những hình bóng ấy rung lên dữ dội rồi tan biến hoàn toàn, không để lại một dấu vết. Không còn tiếng la hét, không còn tiếng bước chân, tựa như ba người họ chưa từng tồn tại.

Im lặng.

Chỉ còn lại ba người họ, những kẻ chưa biến mất, đứng giữa toa tàu rung lắc, giữa ánh đèn nhấp nháy và tiếng còi rú rợn người.

Không ai trong số họ dám ho he một tiếng. Trong thâm tâm, họ đã biết được tại sao lại như vậy.

Nhưng chưa kịp để họ sững sờ lâu hơn, con tàu chao đảo dữ dội, hành lang lắc lư, tiếng kim loại rền vang như sắp vỡ tung. Còi báo động rít lên từng nhịp dồn dập, làm rung cả khối không khí nặng nề, khiến trái tim họ đập loạn xạ theo.

Cậu con trai gằn giọng như muốn xé to bầu không khí chết chóc: ``Không thể đứng im chờ chết\dots\ Phải tìm cách dừng nó lại!''

Trong ánh đèn nhấp nháy đỏ ngầu, mọi thứ trở nên méo mó. Âm thanh hòa trộn thành một mớ hỗn loạn: tiếng còi tàu hú dài, tiếng bánh xe gào thét trên đường ray, tiếng báo động rú inh ỏi, tất cả đập vào tai như một cơn bão sắt thép. Mỗi hơi thở trở nên nặng nề, dồn nén, tựa như cả toa tàu biến thành một nồi áp suất sắp nổ tung.

Hành lang rung bần bật, từng cú chấn động làm ghế ngồi bật gốc, ngã rạp xuống sàn. Cửa kính hai bên vỡ loảng xoảng, những mảnh vụn bắn tung tóe, hắt ánh sáng lóe lên trong từng nhịp đèn chớp. Mỗi bước chân di chuyển đều phải níu chặt vào thanh tay vịn, nếu không sẽ bị quăng quật như món đồ chơi.

Reiko thoáng quay đầu, đôi môi mấp máy định nhắc lại chuyện ``tên thật'' còn bỏ ngỏ. Nhưng chưa kịp thốt thành lời thì Saikyou đã gắt lên, giọng đanh như lưỡi dao: ``\textbf{Không phải lúc này, Reiko!}''

Kyuen liền hét theo, giọng khàn đặc vì gió rít: ``\textbf{Chuyện đó để sau! Giờ ta phải tìm cách nào đó để thoát khỏi đây trước!}''

Reiko nghẹn lời, nước mắt loé lên nhưng cô kịp nén lại, gật mạnh. ``Đ-Được rồi! Vậy phải làm sao?''

Trong đầu của ba người lúc này dấy lên một hình ảnh ghê rợn: con tàu này chẳng khác nào một cỗ quan tài khổng lồ, đang phóng đi với tốc độ siêu thanh. Không còn lối thoát, không còn cơ hội\dots\ chỉ còn một cái chết đang lao tới gần, nhanh đến mức ngay cả ý nghĩ cũng không kịp bám theo.

Saikyou siết chặt thanh vịn, đôi mắt đảo nhanh trong ánh đèn nhấp nháy. ``Cửa thoát hiểm đều khóa kín\dots\ nhưng còn buồng lái!'' Cô chỉ về phía trước, nơi lối hành lang cuối cùng tối om. ``Nếu đến được đó, ta có thể kéo phanh khẩn cấp hoặc đảo cầu điều khiển!''

Tiếng gầm của động cơ như muốn nuốt chửng cả toa. Kyuen hét to trong hỗn loạn: ``Được rồi! Dù chết cũng phải thử!''

Cậu kéo Reiko sát lại bên mình, ánh mắt găm thẳng vào cô em song sinh. ``Em đi trước dẫn đường. Anh sẽ che chắn cho Reiko!''

Saikyou gật mạnh, lao về phía trước, từng bước chân nện xuống sàn thép như búa bổ. ``Giữ chặt, đừng buông tay!''

Reiko run rẩy, nhưng vẫn cố bước theo, giọng cô nhỏ dần, như một lời cầu nguyện lọt thỏm giữa tiếng ồn ào: ``Chỉ cần đến được buồng lái\dots\ chỉ cần dừng được con tàu này\dots''

Ba anh chị em bấu víu lấy nhau, vật lộn giữa cơn địa chấn sắt thép, giữa tiếng còi rú dồn dập như khúc nhạc báo tử. Trong ngực họ, những mảnh ký ức về ``tên thật'' vẫn còn rực cháy, nhưng tất cả bị ép lùi lại phía sau bởi một mục tiêu duy nhất: đến được buồng lái\dots\ hoặc chết trên đường tới đó.

\ct

Khi tới được toa liền kề, cả ba chết sững. Trước mắt họ là một chồng bàn ghế học đường chất ngổn ngang, cao ngất, chặn kín lối đi. Những chiếc bàn gỗ xước xát, ghế nhựa gãy chân, cả bảng đen lẫn tủ sắt nằm chen chúc, như một cơn ác mộng tái hiện khung cảnh lớp học đã bị xé ra và nhồi nhét vào trong toa tàu chật hẹp. Không khí nồng nặc mùi bụi gỗ ẩm và sắt gỉ, bức bối tới ngột ngạt.

Kyuen nghiến răng, lao vào đống đồ, cố hết sức đẩy chúng sang một bên. Nhưng cứ mỗi lần cậu giật mạnh, chúng lại dựng đứng trở lại, ngay lập tức, như có một bàn tay vô hình cố ý sắp đặt.

``Chết tiệt! Nó không nhúc nhích\dots'' Kyuen gầm lên, mồ hôi chảy ròng ròng trên thái dương.

``Không thể nào\dots'' Reiko hốt hoảng, đôi tay nhỏ bé cố nâng chiếc ghế nhưng chỉ càng làm chúng lắc lư dựng dậy.

``Cứ như thể\dots\ có một sức mạnh nào đó can thiệp vậy!'' Saikyou rít lên, vừa kéo vừa đạp, nhưng mọi thứ chỉ càng trở về nguyên trạng, thậm chí còn bừa bộn hơn.

Cả ba người bị dồn đến bước bế tắc. Trong đầu họ lúc này len lỏi một câu hỏi: ``Tại sao chồng bàn ghế này lại xuất hiện ở đây?''

Bất ngờ, tiếng loa rè rè vang lên, khắp toa rung lên trong âm thanh lạo xạo nhiễu sóng. Trong nền tạp âm, một giọng phụ nữ vang vọng, lạnh lùng mà quen thuộc. Nghe như giọng của Shino, nhưng vỡ vụn như vọng từ đáy vực: ``Hãy sử dụng \emph{ánh sáng} để vạch ra lối đi.''

Cả ba người sững người. Kyuen chưa kịp phản ứng thì từ toa trước đó, một vật hình chữ nhật nằm vất vưởng dới sàn, ánh kim loại phản chiếu trong nhịp đèn nhấp nháy. Cậu cúi xuống nhặt lên, đó là chiếc máy ảnh kỹ thuật số.

Cậu liền nhận ra ngay lập tức: vật từng thuộc về Uzuki. Một cô nàng bởm thỏ từng tồn tại, nhưng đã bị vòng lặp nuốt chửng, để lại chỉ là trống rỗng.

``\dots Không còn thời gian nghĩ nữa!''

Cậu đưa chiếc máy ảnh lên, ngón tay run run ấn nút chụp về phía đống bàn ghế chắn lối.

*TÁCH!*

Ánh flash lóe sáng chói lòa, xé toang khoảng tối. Trong tích tắc, chồng bàn ghế tan rã như tro bụi, từng mảnh gỗ, sắt, nhựa hòa vào không khí rồi biến mất, để lại hành lang thông suốt.

Reiko há hốc mồm, gần như quên cả thở: ``Nó\dots\ nó biến mất thật rồi\dots''

Saikyou cũng kinh ngạc, căng tràn cả sợ hãi: ``Hóa ra\dots\ là vậy.''

Họ hiểu ngay: đây chính là công cụ có thể giúp họ sống sót, một lối ra duy nhất khỏi cơn ác mộng đang bủa vây.

Chiếc máy ảnh nằm gọn trong tay Kyuen, nhỏ bé đến mức có thể rơi khỏi tay cậu bất cứ lúc nào. Nhưng với cả ba lúc này, nó chẳng khác gì sợi chỉ mong manh treo lơ lửng trên vực thẳm, một niềm hy vọng nhỏ bé giữa tiếng gầm rú sắt thép đang lao điên loạn.

``Đi tiếp thôi,'' người anh nói, ánh mắt không rời khỏi vật trong tay.

Cả ba lần mò dọc hành lang hẹp, ánh đèn chập chờn như muốn phụt tắt bất cứ lúc nào. Chưa đi được bao xa, trước mặt họ lại hiện ra một chồng bàn ghế khác: cũng lộn xộn, chất cao, chặn kín lối như một bức tường bất khả xâm phạm.

Saikyou nghiến răng, giơ tay định kéo, nhưng Kyuen đã siết chặt vai cô: ``Không cần. Để anh.''

Cậu giơ máy ảnh, chỉnh ống kính hướng thẳng.

*TÁCH!*

Ánh flash lóe sáng, nuốt chửng cả toa tàu trong khoảnh khắc. Chồng bàn ghế rung bần bật như hình ảnh nhiễu sóng rồi tan rã thành từng hạt bụi, biến mất ngay trước mắt. Con đường mở ra.

Reiko thì thầm, run rẩy: ``Nó lại\dots\ biến mất thật rồi. Vậy ra chiếc máy ảnh đó được hoạt động như vậy.

Họ lách qua, và lại một chồng bàn ghế khác chắn ngay phía trước. Kyuen không đợi bàn thêm, lại nhấc máy ảnh lên.

*TÁCH!*

Lại ánh sáng. Lại biến mất.

Cứ như thế, từng bức tường bàn ghế ngổn ngang hiện ra rồi bị ánh đèn nháy cuốn phăng đi, mở ra lối nhỏ dẫn tới cửa toa kế cận. Hơi thở dồn dập, tim đập thình thịch, nhưng không ai dám sải bước quá nhanh. Trong sự nguy cấp, vội vàng chính là cái chết.

Cả ba siết chặt lấy nhau, bước chậm rãi qua cánh cửa sang toa kế tiếp.

\ct

Toa mới mở ra, và họ khựng lại.

Một lần nữa, bàn ghế lớp học chất đống chắn ngang. Những chân ghế chọc ngoằng ngoèo, bàn học nghiêng ngả, trông như một lớp học đã bị xé toạc và nhét vào nơi này.

Saikyou cau mày, buột miệng: ``Tại sao chúng lại ở đây? Bàn ghế lớp học\dots\ trên một con tàu tốc hành? Thật vô lý.''

Reiko ôm chặt lấy cánh tay cô, nhỏ nhẹ đáp: ``Chắc\dots\ chắc là do các vòng lặp chồng lấn. Từng mảnh ký ức, từng mảng không gian\dots\ bị ép vào cùng một chỗ.''

Kyuen im lặng, chỉ đáp lại bằng một cái nhìn, rồi giơ máy ảnh.

*TÁCH!*

Bức tường bàn ghế biến mất. Nhưng ngay khi con đường trống trải, cả ba cùng sững người.

Ở ngay cửa ra vào của toa tàu, có một bóng người.

Đèn nhấp nháy soi mờ gương mặt quen thuộc. Yuka. Cô đứng đó, hai tay điên cuồng đập vào cánh cửa, nhưng tuyệt nhiên không thốt ra một âm thanh nào. Miệng cô mấp máy, nhưng không có lời nói.

Cậu con trai hoảng hốt, gọi to: ``Yuka-san\dots? Tại sao cô ấy lại ở đây?''

Cô em song sinh trố mắt, lắc đầu liên hồi: ``Không thể nào\dots\ Em tưởng cô ấy đã biến mất rồi mà\dots''

Cô em út lùi lại một bước, nhìn về phía hai anh chị, lên tiếng ngăn cản: ``Đừng. Đừng lại gần.''

Nhưng cả ba vừa bước thêm nửa nhịp, hình ảnh cô nàng tai mèo bỗng rung lên dữ dội như một đoạn video lỗi khung hình, rồi tan biến vào hư không, để lại khoảng trống lạnh ngắt.

Một ảo ảnh. Rõ ràng là một ảo ảnh.

Mồ hôi túa ra trên đầu họ. Một cảm giác lạnh sống lưng. Cả ba đều hiểu: đó không còn là Yuka thật, mà chỉ là một vết nhòe còn sót lại của vòng lặp.

Kyuen thấy ngực mình như bị bóp nghẹt: hình ảnh Yuka vừa tan biến khiến cậu nhớ đến Kaoru - cũng cái cách biến mất đầy đột ngột ấy, cũng cái khoảnh khắc không kịp trăn trối.

Saikyou thì sững sờ tới độ không nói lên lời. Cô không khóc, chỉ thấy lòng bàn tay nóng rát vì móng đã đào sâu vào thịt. Với cô, ảo ảnh Yuka là lời cảnh tỉnh phũ phàng: nếu không đủ mạnh, cô sẽ biến thành ``vệt nhòe'' tiếp theo. Quyết tâm của cô lên ngọn lửa, phải tìm ra thoát khỏi nơi này trước khi mọi chuyện trở nên tệ hơn.

Reiko run nhè nhẹ, hai vai thu lại như con chim nhỏ trước cơn bão. Cô là người duy nhất không thể che giấu nỗi sợ: mỗi lần một bóng người tan biến, cô lại thấy mình nhỏ bé hơn. Nhưng khi nhìn thấy ánh mắt của người anh lớn - đỏ ngầu nhưng kiên định - cô bất giác siết chặt tay áo cậu, thầm nhủ: ``Mình không được trở thành gánh nặng.'' Nỗi sợ không biến mất, nhưng trong mắt cô bắt đầu lấp lánh một mầm quyết tâm mong manh: phải sống, phải nhớ, và phải cứu lấy chính mình.

Họ tiếp tục đi, thêm hai lần chồng bàn ghế chắn ngang.

*TÁCH!*

*TÁCH!*

Ánh flash xóa tan mọi vật cản, lối đi thông ra.

Cuối cùng, trước mắt họ là cánh cửa dẫn sang toa kế cận. Họ dừng lại, cả ba thở hổn hển để lấy lại tinh thần, mồ hôi nhễ nhại, nhưng ánh mắt lại ánh lên quyết tâm.

``Lên nào.''

\ct

Toa kế tiếp mở ra, và ngay lập tức cả ba phải dừng bước. Bàn ghế lớp học chất cao, chồng chéo lộn xộn, chắn kín gần như toàn bộ lối đi. Kyuen giơ máy ảnh, bấm một cú chớp.

*TÁCH!*

Ánh flash nuốt chửng toa tàu trong một nhịp. Bàn ghế run rẩy, méo mó rồi tan biến\dots\ nhưng lần này, thay vì biến mất hoàn toàn, chồng bàn ghế chỉ tan rã một phần. Vẫn còn sót lại một chiếc bàn chắn ngang, cao vừa đủ để buộc cả ba phải trèo qua.

Kyuen bước đến đầu tiên, chân sải một bước, gọn gàng vượt qua. Áo khoác phất nhẹ theo động tác, hoàn toàn không gặp vấn đề gì. Nhưng khi quay lại, cậu thoáng khựng lại, nhìn hai cô em gái đang nhìn chiếc bàn với tâm trạng có phần đối nghịch: một ngán ngẩm, một e dè.

``\dots À, phải. Anh quên mất là hai đứa đang mặc váy.''

Saikyou nhếch môi, chẳng buồn do dự. Cô chống tay lên mép bàn, thoáng nghiêng người nhảy qua một cách mạnh bạo. Váy ngắn xoè ra trong thoáng chốc, nhưng cô chẳng hề bận tâm, chỉ liếc sang ông anh song sinh một cái đầy thách thức: ``Có gì mà phải ngại, đúng không?''

``Ít nhất cũng phải biết xấu hổ chứ, đồ ngốc\dots'' người anh song sinh đảo mắt, lẩm bẩm.

Reiko thì khác. Cô đứng chết trân, mắt bối rối nhìn hai anh chị, tay giữ chặt gấu váy, gương mặt đỏ ửng. ``Em\dots\ Em không thể\dots\'' giọng cô lí nhí, run run.

Cô chị quay lại, cau mày: ``Reiko, giờ không phải lúc để ngại ngùng. Chúng ta đang chạy đua với tử thần đấy.''

``Nhưng mà\dots\ nếu có người nhìn thấy\dots''

``Không ai nhìn đâu, cả tàu chết lặng rồi.'' Kyuen nhẹ giọng, vẫn quay lưng. ``Anh đếm: một, hai, ba.''

Cô mím chặt môi, nước mắt lấp lánh nơi khóe mi, nhưng rồi hít một hơi thật sâu. Cô xoay người, quỳ gối lên bàn, cố gắng giữ váy che chắn hết mức có thể, rồi chậm rãi trèo qua. Cậu con trai lặng lẽ đưa tay ra đỡ, không một lời.

``Cảm\dots\ cảm ơn anh, Onii-sama\dots'' cô thì thào, mắt nhìn xuống.

Họ tiếp tục, lại thêm một chồng bàn ghế nữa. Lần này, ánh flash xóa sạch, không để lại chút gì.

Nhưng khi tới cuối toa, cả ba người bỗng chốc dừng lại.

Ở những hàng ghế phía xa, trong ánh đèn nhấp nháy mờ ảo, có ba bóng đen lấp ló. Một mang dáng tomboy tóc ngắn, vai xuôi mạnh mẽ; không ai khác ngoài Kaoru.

Một dáng khác nhỏ nhắn với chỏm tóc ngố quen thuộc, là Saya. Và một dáng cao gầy với mái tóc buộc đuôi ngựa dài chấm lưng, tai sói lộ rõ, đó là Miku.

Không ai trong ba anh chị em dám cất lời. Không khí như đông cứng, nhịp tim dồn dập át cả tiếng kim loại nghiến ken két dưới bánh tàu.

Reiko nắm chặt lấy tay áo người anh, đôi mắt long lanh. Cậu siết tay cô em, không nói gì, chỉ dán mắt vào ba cái bóng.

``Họ\dots\ cũng từng là bạn của chúng ta,'' cô thì thào, giọng như sắp vỡ. ``Dù chỉ là một khoảng thời gian ngắn.''

``Giờ chỉ còn là bóng ma của vòng lặp,'' Kyuen đáp, giọng khàn đặc. ``Đừng chạm vào nhé, hai đứa.''

Saikyou siết chặt nắm tay, mắt không chớp. ``Nếu họ xuất hiện ở đây, tức là chúng ta đang đi đúng hướng.''

\dots

\dots

Bất ngờ, chúng giật nảy lên như hình ảnh bị nhiễu. Sau vài nhịp co giật, cả ba ảo ảnh tan biến trong không khí, để lại khoảng trống lạnh ngắt.

Không ai nói gì thêm. Họ lẳng lặng tiến lên, mở ra cánh cửa dẫn tới toa tiếp theo.

\ct

Toa mới lại không khác mấy: bàn ghế chất đống chặn ngang. Mỗi chồng mỗi kiểu: có chồng biến mất hoàn toàn ngay khi ánh flash lóe sáng, có chồng chỉ vỡ vụn một nửa, buộc họ phải lách người, hoặc gồng mình trèo qua.

Kyuen liên tục bấm máy, ánh sáng lập lòe từng hồi, soi rõ khuôn mặt căng thẳng của cả ba.

Và rồi, khi một đống bàn ghế vừa tan biến, bất chợt một bóng người vụt hiện ra ngay trước mặt họ.

``Akane\dots''

Cả ba sững lại. Đó chính là cô gái tóc hai bím đuôi sam quen thuộc, ánh mắt hốt hoảng, đôi tay vung vẩy. Cô chạy vụt về phía cuối toa, vừa chạy vừa hét lên, giọng lạc đi vì khẩn thiết: ``\textbf{Yuka! \dots Yuka!! Bà ở đâu!?}''

Ba anh chị em chỉ kịp nhìn theo, không ai kịp đưa tay ra. Saikyou vô thức giơ tay như muốn gọi lại theo thói quen, nhưng Akane không quay đầu. Cô chỉ chạy, chạy mãi về phía cánh cửa toa kế tiếp, cho đến khi thân ảnh mờ đi, nhòe nhoẹt như một đoạn phim lỗi khung hình, rồi tan biến.

Im lặng nặng nề đổ ập xuống.

Reiko run rẩy, giọng nghẹn lại: ``\dots Họ\dots\ Quả thật họ đều từng tồn tại thật mà.''

Kyuen cúi mặt, siết chặt chiếc máy ảnh trong tay, trong lòng dâng lên cảm giác đau nhói. Saikyou siết chặt nắm tay, mím môi, nhưng không đáp. Chỉ có tiếng tim cô đập gấp gáp át lên cả tiếng rít của bánh sắt.

Không ai nói thêm lời nào. Chỉ còn tiếng bánh tàu nghiến ken két, tiếng còi hú rít dài. Cùng với đó là những ánh mắt đan xen sự mệt mỏi, quyết tâm, và nỗi đau chưa kịp gọi tên.

\ct

Cánh cửa toa khẽ mở, kêu lên một tiếng két dài rồi lộ ra trước mắt họ một chồng bàn ghế ngổn ngang, chất đống lộn xộn y hệt những toa trước.

Kyuen nhấc chiếc máy ảnh trong tay, nhắm thẳng vào chồng bàn ghế và bấm chụp.

*TÁCH!*

Ánh flash sáng loé khắp toa, nhưng đống bàn ghế trước mặt vẫn sừng sững, không hề biến mất.

Cậu nhíu mày, thử lại lần nữa.

*TÁCH!*

\dots

Rồi lại thêm một lần nữa.

*TÁCH!*

\dots

Ba lần liên tiếp, nhưng chẳng có gì thay đổi. Chồng bàn ghế vẫn ở đó.

Kyuen lẩm bẩm, giọng mất kiên nhẫn: ``Sao lại không được nhỉ\dots?''

Saikyou chau mày, lùi lại nửa bước: ``Đừng bảo là cái máy đã\dots\ hết tác dụng rồi nhé?''

Reiko khẽ lắc đầu, nhưng giọng cô đầy căng thẳng: ``Không\dots\ em nghĩ là càng đi xa, chiếc máy ảnh này càng mất dần sức mạnh. Có lẽ nó vốn không phải để dùng mãi mãi\dots''

Đúng lúc ấy, một bóng đen hiện ra ngay trước chồng bàn ghế. Bóng đen ấy mang hình dáng của Uzuki. Hai tay nó đưa ra, run rẩy như đang ``xin'' lại chiếc máy ảnh.

Cả ba người đều giật mình.

Cô chị lớn thốt lên, lùi thêm một bước: ``Ơ\dots\ Uzuki-san\dots?

Kyuen siết chặt máy ảnh, mồ hôi túa ra hai bên thái dương. Nhưng rồi cậu hít một hơi sâu, đánh liều: ``Nếu đây\dots\ là cách duy nhất\dots\ thì thử vậy!''

Cậu chìa máy ảnh ra. Bóng đen Uzuki từ tốn đón lấy, và ngay khi ngón tay nó chạm vào nút chụp\dots

*CHOÉ!*

Một luồng ánh sáng cực mạnh nổ tung khắp toa, làm cả ba phải nheo mắt, che mặt. Khi ánh sáng mờ dần đi, đống bàn ghế trước mặt đã biến mất. Nhưng đồng thời, cả chiếc máy ảnh lẫn bóng đen Uzuki cũng không còn nữa.

Saikyou thì thào, giọng còn run: ``Vừa rồi\dots\ là gì vậy?''

Reiko nhíu mày, không giấu được vẻ hoang mang: ``Em\dots\ em không biết\dots''

Kyuen lấy chiếc đèn pin ra, thở dài: ``Dù gì thì đường cũng thông thoáng rồi. Đó\dots\ cũng là ân huệ cuối cùng mà cậu ấy làm cho chúng ta.''

Saikyou cắn môi, nhìn về phía xa: ``Giờ chỉ còn cách tiến về phía trước thôi.''

Cô em út gật nhẹ, giọng run run: ``Đúng\dots\ không thể quay đầu lại được nữa rồi.''

May mắn thay, trước mắt họ bây giờ không còn vật cản nào nữa. Đường đi thẳng tắp, trống trải, và dường như tiếng đầu máy kéo vọng lại càng lúc càng rõ, càng lớn. Niềm tin rằng họ đã rất gần với buồng lái len lỏi trong cả ba người.

Trước mặt họ bây giờ là một cánh cửa toa khác. Tiếng lục cục của đầu kéo xe lửa lớn tới mức khiến họ cũng phải nhíu mày.

Saikyou bước lên, đặt tay vào nắm cửa cuối toa và cố đẩy\dots\ nhưng cánh cửa không nhúc nhích. Cô thử xoay, thử kéo, nhưng vẫn chẳng mở được.

Cô nghiến răng, giọng căng cứng: ``Cửa\dots\ bị khóa rồi.

Kyuen tiến lên, thử đẩy mạnh: ``Không có dấu hiệu gì là mở được cả\dots''

Reiko cầm chiếc đèn pin từ anh trai, rọi vào ổ khóa trên đó: ``Có lẽ\dots\ cần chìa khóa chăng?''

Ngay lúc đó, toàn bộ bóng đèn trên tàu tắt phụt, cả đèn pin cũng không còn hoạt động, để lại một màu đen đặc nghẹt. Dù vậy, đoàn tàu vẫn lao đi ầm ầm, như thể nó chẳng hề quan tâm gì tới số mệnh của chính nó.

Cô em út rùng mình, giọng ren rén: ``Đèn\dots\ tắt hết rồi\dots''

Saikyou siết chặt tay, giọng trầm: ``Đừng tách ra\dots\ dù có chuyện gì.''

Từ đầu toa, một luồng sáng yếu ớt bất ngờ rọi xuống. Trên một chiếc ghế, có bóng dáng một người đang ngồi.

Reiko mặt tái mét, thì thào: ``Có\dots\ có người ở đó\dots''

*PHỤT!* Ánh sáng vụt tắt.

Rồi bật trở lại, chiếc ghế kia giờ đã tiến sát hơn về phía họ. Trong khoảng sáng mờ, họ thấy thấp thoáng một mái tóc màu tím quen thuộc.

Saikyou nín thở: ``Không lẽ\dots\ Yuka-san\dots?''

*PHỤT!* Đèn lại tắt.

Khi sáng lên lần nữa, chiếc ghế đã gần sát hơn nữa, chỉ còn cách họ vài mét. Giờ đây, họ nhìn rõ đó chính là Yuka, ngồi bất động, mắt cúi xuống.

Kyuen lẩm bẩm: ``Yuka-san\dots\ lại nữa ư?''

*PHỤT!* Bóng tối lại bao trùm.

Cả ba người nín thở, tim đập loạn nhịp. Rồi, ánh sáng trở lại, lần này chiếc đèn pin cũng hoạt động. Chiếc ghế ngay sát trước mặt họ, và Yuka đang ngồi đó, bất động, không phản ứng.

Saikyou bước chậm, giọng khẽ: ``Cô ấy\dots\ tại sao lại ở đây vậy?''

Reiko nắm chặt tay áo, giọng run: ``C-Cẩn thận đấy, Onee-sama\dots''

Họ từ từ tiến lại gần. Và ngay khoảnh khắc chỉ còn cách một bước chân, ảo ảnh Yuka vụt tan biến như khói. Trên ghế, chỉ còn lại một chiếc chìa khóa vàng lạnh lẽo.

Cậu con trai cúi xuống, nhặt nó lên bằng hai ngón tay, giọng khàn đặc: ``Cái này\dots\ không chỉ là chìa khóa.'' Cậu xoay nhẹ, ánh đèn lướt qua bề mặt, loé lên tia sáng xanh lạnh. ``Nó nặng hơn nó vẻ. Như thể\dots\ ai đó đã nhét cả một mảnh ký ức vào trong đó.''

Saikyou nhìn về phía cửa cuối toa, môi mím chặt: ``Cánh cửa đó vừa xuất hiện. Không có tay nắm, không có khe hở, chỉ có lỗ khóa.'' Cô đưa tay lên vách, khẽ gõ, nghe tiếng kim loại vang lên rỗng tuếch. ``Chúng ta không được chọn nữa.''

Reiko ôm lấy cánh tay mình, giọng run nhè nhẹ: ``Em\dots\ nghe thấy tiếng tim mình đập. Nhưng mà\dots\ tiếng tim của ai?'' Cô nhìn chằm chằm chiếc chìa, mắt ánh lên nỗi sợ không tên. ``Nếu mở ra\dots\ liệu chúng ta còn là chúng ta không?''

``Không ai trả lời được,'' Kyuen thì thầm, tiến về phía cửa. ``Nhưng ở lại đây thì chắc chắn sẽ mất.'' Cậu đưa chìa vào ổ, xoay nhẹ. Một tiếng *tách* khẽ như xương rời khớp. Cánh cửa nặng nề tách ra, khói nóng và mùi dầu cháy xông vào mặt.

Trước mắt họ là buồng lái chật hẹp, như một cái bụng sắt đang co rút. Tay lái rung lắc dữ dội, đồng hồ tốc độ kim đã chạm vạch đỏ, rung như muốn bật ra ngoài. Dây cáp lòng thòng, van xoay hoen gỉ, ánh đèn vàng mờ như đang thoi thóp hơi thở cuối.

Saikyou bước vào trước, giọng gằn qua tiếng rít kim loại: ``Phanh chính bị khóa cứng. Tìm van phụ\dots\ có lẽ nằm dưới sàn!'' Cô quỳ xuống, tay sờ soạng dưới bệ lái, mồ hôi lạnh chảy dọc thái dương.

Kyuen lao tới cần gạt lớn, hai tay giật mạnh, gân xanh nổi lên: ``Không nhúc nhích\dots\ Chặt quá!'' Cậu quay sang Reiko, mắt đỏ ngầu. ``Em có thấy nút đỏ nào không? Cái loại `dừng khẩn cấp' ấy!''

Reiko lục lọi bảng điều khiển, giọng gấp gáp: ``Dưới đây! Có kính vỡ rồi\dots'' Cô dùng khuỷu tay đập mạnh, mảnh kính cắt vào da, máu rỉ ra.

``Ái-da\dots''

``Reiko!'' cô chị trố mắt ra, nhìn khuỷu tay đang rỉ máu.

``Em\dots\ Em ổn,'' cô em út đáp lại.

Cô ấn nút đỏ, cả buồng lái rung chuyển, một tiếng *RẦM* nổi lên từ gầm tàu.

Cậu con trai gầm lên: ``\textbf{Giữ chặt vào! Nó đang lật đấy!}'' Cậu vươn tay, níu lấy dây kéo trên trần, kéo hai em sát vào tường. ``\textbf{Đừng buông ra nha, hai đứa!}''

Con tàu gào thét tiếng kim loại, nghiêng hẳn. Đèn vỡ tung, đồng hồ tốc độ bật ra, bay qua đầu họ. Sàn nghiêng đến 45 độ, ba thân người trượt xuống, tay cào vào kim loại rít lên những vệt máu loang lổ. Tia lửa đỏ rực tóe ra ngoài cửa sổ, như máu của chính con tàu.

Reiko khóc nấc, giọng nhỏ xíu: ``Chúng ta\dots\ sắp chết rồi sao?'' Cô siết chặt tay Saikyou, móng tay cắm vào da chị.

Kyuen cắn răng, máu me bên khóe môi: ``Không\dots\ chưa phải lúc\dots'' Cậu cố trườn lên, tay vớ lấy chiếc đồng hồ đeo tay đã lắp ráp từ trước, mặt kính đã nứt nẻ sau cú va chạm.

Cậu ấn mạnh nút bên hông, chiếc đồng hồ rung lên nhè nhẹ, như mạch máu đầu tiên được nối lại. Mặt đồng hồ sáng lên tia trắng nhỏ, lan từ tâm ra ngoài, rồi bùng lên thành luồng sáng xoáy tròn.

Cô em gái song sinh nhìn chằm chằm ánh sáng, mắt trợn to: ``Cái gì thế này!?''

Reiko nhỏ giọng như từ trong mơ: ``Ánh sáng\dots\ ấm quá\dots''

Luồng sáng trắng bao trùm cả buồng lái, xuyên qua thân thể họ, làm mờ mọi đường nét. Tiếng kim loại rít lên chậm dần, tiếng còi tàu xa dần, như thể thế giới đang xa dần trong tầm tay.

Kyuen mỉm cười, máu chảy dọc khóe môi: ``Thì ra\dots\ thời gian cũng sợ\dots\ bị bỏ lại\dots''

Saikyou buông lỏng tay, thì thầm: ``Ta\dots\ sắp gặp lại mình\dots\ của một kiếp khác\dots''

Reiko chớp mắt lần cuối, nước mắt lăn ngược vào thái dương: ``Em\dots\ nghe thấy tiếng ai gọi tên\dots\ nhưng\dots\ không phải tên mình\dots''

Ánh sáng bùng lên rực rỡ, nuốt trọn cả đoàn tàu. Ba thân người nhẹ bẫng, trôi trong khoảng trắng không trọng lượng. Rồi tất cả im bặt, không tiếng, không hơi thở, không còn gì ngoài màu trắng xoá\dots

\dots

\dots

\ct

\pov{???}

\dots

\dots

Một giọng nói nữ vọng lại từ đâu đó trong không gian mờ ảo.

``Các cậu thấy trường học có vui không?''

Giọng nói ấy nghe như đang vang vọng từ rất xa, nhưng đồng thời cũng như thể đang thì thầm ngay bên tai.

\dots

``Tớ dám chắc là các cậu thấy vui.''

Âm sắc trong giọng nói ấy thật dịu dàng, ấm áp\dots\ và kỳ lạ thay, nó khiến tôi cảm thấy nhói đau.

BÍNH BONG\dots\ BÍNH BONG\dots

Tiếng chuông báo hết giờ vang vọng, như một tiếng vọng xa xăm từ quá khứ.

\dots

``\dots Đã hết tiết rồi đó. Thôi nào, ta hãy cùng về thôi.''

\dots

Giọng nói ấy lại cất lên. Một giọng nói quen thuộc đến kỳ lạ.

Tôi biết giọng này. Tôi chắc chắn là mình biết.

\dots

\dots

Sao tôi cảm thấy\dots\ điều này quen thuộc vậy?

Tôi biết rất rõ đó là giọng của Shino-san. Không thể lầm vào đâu được.

Nhưng mà\dots\ tại sao cô ấy lại ở đây?

\dots

\dots

Tôi giật mình mở mắt. Cả người ê ẩm, đầu vẫn còn ong ong như sau cú va đập mạnh. Phải mất một lúc tôi mới lồm cồm bò dậy, tay chống xuống nền lạnh buốt.

Tôi đang ở đâu đây?

Trước mặt tôi là một hành lang trường học. Cửa sổ cao, những bóng đèn sáng chập chờn. Ánh trăng hắt vào, kéo bóng tôi dài ra méo mó trên nền gạch. Nơi này vừa xa lạ\dots\ vừa quen thuộc đến rợn người.

``Kuon-user.''

Tôi khựng lại. Một giọng đàn ông khản đặc, pha chút âm sắc Pháp, vang lên đâu đó ngay sát bên.

``\eng{Kuon-user, cậu tỉnh lại rồi chứ?}''

\dots Tiếng Anh? Ông ta\dots\ nói chuyện với tôi bằng tiếng Anh ư? Tôi không có vấn đề gì về trình độ ngoại ngữ của tôi, nhưng\dots\ tại sao lại là lúc này?

\dots Khoan đã, tôi hiểu được ngoại ngữ từ khi nào vậy?

Quan trọng hơn hết, ai đang nói chuyện với tôi?

``\eng{Kuon-user, cậu ổn không?}''

Tôi nhìn quanh. Không một bóng người. Nhưng ngay trước mặt, trên sàn, một chiếc đồng hồ điện tử với màn hình lập lòe sáng yếu ớt. Chính nó vừa cất tiếng.

Tôi nuốt khan. ``\eng{\dots Tôi á? Ý ông gọi ai cơ?}''

``\eng{Cậu chứ còn ai nữa, Kuon-user. Không cậu thì ai?}'' giọng ông ta cất lên, có chút gì đó bực dọc.

Tôi cứng người. Một cơn ớn lạnh chạy dọc sống lưng. Tôi? Không phải tôi là Haragen Kyuen sao?

Tôi ngồi phịch xuống, cố ép trí óc nhớ lại. Hình ảnh tàu điện lật nghiêng, rồi màn đêm nuốt chửng\dots\ Tôi lắc mạnh đầu, nhưng tất cả vẫn mờ mịt.

``Chờ đã\dots\ `Tôi'? Cái tên đó\dots\ giống như trong\dots\ kỷ yếu\dots''

\dots Phải. Cái kỷ yếu mà ba anh chị em chúng tôi, cùng Uzuki-san và Sora-san - mà thực chất là bản ngã của Shino-san - đã đọc. Những cái tên xa lạ mà thân thuộc ấy\dots\ xuất hiện: Kuon. Kaigou. Suzuki.

Tôi cắn chặt răng, ngực nặng như bị đè ép. Vì sao đồng hồ này lại gọi tôi như vậy? Và vì sao nó khiến tôi cảm giác như mình từng nghe thấy, từng sống với cái tên đó?

``\eng{Tôi biết là cậu vừa tỉnh lại, nên cứ bình tĩnh lại nào, Kuon-user. Cậu không cần phải gắng quá đâu.}''

Tôi vịn vào tường, hít thở gấp. Trong đầu vẫn quanh quẩn cái tên ấy: Kuon. Nó vừa lạ lẫm, vừa quen thuộc đến rùng mình.

Ánh trăng chiếu xuống, lướt qua những dãy bàn học cũ kỹ đặt ngổn ngang dọc hành lang. Trong lúc bước đi loạng choạng, tôi thoáng thấy một vật nhỏ phản chiếu ánh sáng trên mặt bàn.

Một tấm thẻ nhựa.

Tôi khựng lại, tim đập mạnh. Trên nền thẻ màu lam bạc, dòng chữ hiện rõ:

``Trường Đại học Xây dựng Kawachi.

Khoa Công nghệ thông tin.

Hikawa Kuon (火川久遠).'' Đi kèm với đó là ảnh chân dung, có nét mặt rất giống tôi.

\dots Không. Đó chính là tôi.

``\dots!?''

*KENG!*

Trong thoáng chốc, đầu tôi như bị ai đó dùng búa tạ giáng xuống. một cú va đập chát chúa ngay giữa trán, khiến toàn bộ hộp sọ rung lên bần bật. Cơn đau nhói buốt lan tỏa tức thì, từng mạch máu như bị thắt nút, từng dây thần kinh như bị kéo căng đến đứt phựt. Mắt tôi nổ đom đóm sáng trắng, tai ù đặc, không khí xung quanh bỗng nặng trĩu như chì.

``aaaaa\textbf{AAAAhhhhHHH!!}''

``Kuon-user!?''

Tôi ôm lấy thái dương, ngón tay siết chặt đến mức muốn bấm sâu vào da thịt, như thể chỉ cần ép mạnh hơn thì có thể bóp nghẹt cơn đau đang bùng phát bên trong.

\hll.

Tôi thấy mình đang ngồi giữa văn phòng, tập giấy tờ chất cao như núi. Những gương mặt quen thuộc, như Towa-chan và Ru-chan rụt rè hỏi tôi nên làm gì để được là chính mình; A-chan, bình thản gõ bàn phím, cùng tôi giải quyết mớ sổ sách ngổn ngang.

Những tiếng cười, những buổi họp, những lần tôi đứng ra giải quyết mớ hỗn độn mà các cô gái gây ra. Tất cả ùa về, vừa thân thuộc vừa nhói đau.

Rồi khung cảnh thay đổi. Một trạm xe buýt, bầu trời sáng cuối thu se lạnh. Tôi đứng đó, và rồi trong lúc mải đuổi theo đồng tiền xu, tôi bị một vòng xoáy - kết quả từ những con rối của Shion-chan - gây nên. The Void.

Trong bóng tối đó, hai gương mặt: một cô gái tóc đuôi ngựa, một cô gái tóc cột hai bên đi cùng với tôi, sau khi tôi giúp họ đoàn tụ.

Hikawa Kaigou-chan\dots\ và Hikawa Suzuki.

Tiếp đó, một thế giới kỳ lạ: ``There Is No Game.'' Ba chúng tôi, cùng nhau, và rồi\dots\ một lỗ xoáy khác xé toạc mọi thứ. Nijisanji. Mười hai ngày. Sau đó\dots\ trắng xóa. Ký ức dừng lại ở đó.

Tôi thở hổn hển, mồ hôi lạnh túa ra trên trán.

``\eng{Kuon-user? Cậu có ổn không vậy?}''

Giọng nói khàn khàn của chiếc đồng hồ lại vang lên, lần này nghe gần gũi hơn.

``\eng{Ừ\dots\ Tôi vẫn ổn. Chỉ là hơi chập choạng thôi\dots\ Game-san.}''

Tôi siết chặt tấm thẻ sinh viên trong tay, trái tim đập thình thịch.

Không phải Haragen Kyuen.

Tôi là \emph{Hikawa Kuon}. Đây mới là tôi.

Tôi ngẩng đầu lên, nhìn hành lang tối om trải dài trước mặt. Trong bóng tối, tôi bỗng thấy mình không còn hoàn toàn đơn độc.

Tôi khẽ lắc đầu, cố gắng xua đi cơn đau còn sót lại. Bàn tay vẫn giữ chặt tấm thẻ sinh viên, như thể nó là sợi dây duy nhất neo giữ tôi vào thực tại.

Tôi từ tốn nhặt chiếc đồng hồ đeo tay lên, quàng vào cổ tay trái, giống như mọi khi. Thật kỳ lạ, cảm tưởng như tôi đã quên mất một việc gì đó rất quen thuộc từ lâu, mà giờ cuối cùng cũng nhớ lại.

``Game-san này\dots'' tôi gọi khẽ, giọng khàn đặc, ``\eng{ông biết tại sao tôi lại lạc vào nơi này không?}''

Chiếc đồng hồ trên cổ tay khẽ rung, rồi vang lên giọng trầm, khản, có chút mệt mỏi như thể chính nó cũng vừa tỉnh dậy sau một giấc ngủ dài. ``\eng{Không, Kuon-user. Tôi không biết. Đến tôi còn không hiểu chuyện gì đã xảy ra với tôi nữa. Chỉ có một điều chắc chắn: nơi này\dots\ không phải thế giới mà cậu cùng hai người kia muốn tới.}''

Tôi nhíu mày. ``\eng{Ông chắc chứ?}''

``\eng{Chắc. Tôi\dots\ mù tịt về không gian này. Và chính điều đó mới khiến tôi thấy\dots\ lạ.}''

Tôi nuốt khan, đưa mắt nhìn quanh.

Dưới ánh sáng huỳnh quang mờ mịt, một cánh cửa trượt phía xa phản chiếu thứ ánh sáng yếu ớt như dẫn lối. Một cảm giác quen thuộc bất chợt dâng lên, khiến sống lưng tôi lạnh toát.

``\eng{Đ-Đây là\dots}''

Ký ức từ ``Haragen Kyuen'' - cái tên giả mà tôi từng sống cùng - dội về như những mảnh ghép bị vỡ. Những hành lang dài, những lớp học giống hệt nhau\dots\ Đây chính là hành lang trường Aogami. Nhưng khác biệt là: mọi thứ hiện tại mang vẻ ngoài hiện đại, không phải khung cảnh xưa cũ trong trí nhớ.

Tôi bước chậm rãi dọc hành lang. Một bên là những ô cửa sổ khổng lồ, ngoài kia chỉ đặc quánh một màu đen vô tận. Phía đối diện là dãy lớp học, mỗi lớp đều có cửa trượt đôi và tấm bảng lục dán thông báo. Không có số lớp. Không có bất kỳ phân biệt nào.

``\eng{Thật\dots\ kỳ lạ,}'' tôi khẽ lẩm bẩm.

``\eng{Mọi lớp học giống hệt nhau,}'' ông ta đáp lại. ``\eng{Dữ liệu này không mang tính thực tế. Trông nó giống như một mô hình được dựng qua trí nhớ ai đó.}''

Tôi khẽ rùng mình. Tôi ngó mắt nhìn vào một lớp gần đó.

Bên trong là những gì quen thuộc của một lớp học bình thường: bảng phấn xanh, bốn dãy bàn ghế thẳng hàng, một chiếc máy chiếu treo trên trần. Đèn huỳnh quang trắng nhợt chiếu xuống\dots\ và trên tường, chiếc đồng hồ cơ đứng im lìm ở một mốc duy nhất.

10:10.

Tôi run lên, thì thầm gần như không thành tiếng:
``\dots Không thể nào\dots''

Game-san im lặng một lúc, rồi cất lời, thấp và nặng nề:
``\eng{10:10\dots\ Kuon-user, có vẻ\dots\ đây không phải một con số ngẫu nhiên. Cậu biết gì về nó không?}''

``\eng{\dots Có, tôi biết. Rất rõ là đằng khác.}''

Ánh mắt tôi bỗng tối sầm khi ký ức đâm xuyên vào tâm trí: mảnh đất sạt lở, tiếng gào thét, và ngôi trường sụp đổ chỉ trong khoảnh khắc.

Một hơi thở dài bật ra, lẫn trong cơn run nhẹ. Tôi bắt đầu nói, như thể chính mình đang lần ngược lại cuốn nhật ký của một kẻ xa lạ, nhưng cũng chính là tôi.

``\eng{Tất cả bắt nguồn từ\dots\ `Nếu không đối xử tốt với bạn chuyển tới, một sự kiện kinh hoàng sẽ giáng xuống ngôi trường Aogami.' Lời nguyền đó\dots\ ban đầu chỉ như một trò đùa truyền miệng. Nhưng tôi biết, nó bắt nguồn từ Shino-san. Không đơn giản như người ta tưởng.}''

Tôi ngẩng nhìn trần hành lang, ánh đèn huỳnh quang phản chiếu trên khóe mắt.

``\eng{Mọi thứ phức tạp hơn vậy rất nhiều. Shino mang hai bản ngã. Một bản ngã cứ mãi chạy trốn thực tại. Một bản ngã khác thì chấp nhận nó, muốn người kia cũng làm vậy. Và cuối cùng, cô ấy đã thành công thuyết phục nhân cách chối bỏ hiện tại kia quay trở về. Nhưng ngoài ra\dots\ tôi đã nhìn thấy tận mắt những ảo ảnh vỡ vụn. Những cánh hoa bỉ ngạn đỏ mọc lên từ nhà ga\dots\ những dấu hiệu bất thường mà không một ai khác để ý.}''

Game-san im lặng một lúc, rồi cất giọng trầm ấm nhưng đầy cân nhắc: ``\eng{Nếu Shino thật sự mang hai bản ngã, thì cả hai đều đang cùng tiến tới một điểm đích, nhưng theo hai hướng khác nhau. Nói theo cách dân gian thì: một người cứ mãi quanh quẩn trong vòng luẩn quẩn, một người thì dừng lại để nhìn nhận sự thật. Cánh hoa bỉ ngạn đỏ\dots\ có thể là mối nối giữa hai thế giới, những vết nứt mà chỉ người từng chịu tổn thương mới thấy. Cậu đang đứng giữa ranh giới của ký ức và hiện thực, Kuon-user.}''

Tôi khựng lại, mắt mở to. ``\eng{Vậy là\dots\ những gì tôi thấy chỉ là lời thủ thỉ của một ai đó?}''

``\eng{Không chỉ thế,}'' ông thì thầm, ``\eng{nếu nơi này đã chia đôi Shino, thì cũng có thể đã chia đôi cậu. Có lẽ `Kyuen' chỉ là cái tên mượn, còn `Hikawa Kuon' mới là tên thật mà nơi này đang cố giấu. Cậu đã sống trong một phiên bản được tạo ra, còn bản gốc\dots\ có thể đã bị ẩn từ lâu.}''

``\eng{Và đích thân tôi\dots\ là người đã hiện nó?}''

``\eng{Chính xác.}''

``\textit{Với tấm thẻ học sinh ư!?}'' tôi muốn nói ra, nhưng giữ lại trong lòng. Dù biết kịch bản này rất khó tin, nhưng không ngờ nó lại xảy ra thật.

Tôi cười nhạt, có phần chua chát. Tiếng bước chân của tôi vang khẽ trên nền gạch xanh bóng loáng.

``\eng{Nhưng điều kỳ lạ nhất\dots''} tôi tiếp tục, lang thang dọc theo hành lang, ``\eng{là những bạn học mà tôi từng gặp, trong suốt thời gian sống dưới cái tên Kyuen. Uzuki-san, Akane-san, Yuka-san\dots\ trông chẳng khác gì Peko-chan, Koro-chan và Okayu-chan. Chỉ khác là không có tai, không có đuôi. Rồi Nanase-san, Yae-san, Hanabi-san, Saki-san chính là hình ảnh phản chiếu của Sui-chan, Noel-PON, Mat-chan và Aki-chan. Tôi\dots\ không thể coi đó chỉ là trùng hợp nữa.}''

Chiếc đồng hồ khẽ rung lên, rồi giọng Game-san vang ra, có chút trầm ngâm: ``\eng{\dots Một thế giới song song, Kuon-user. Ở đó, những cái tên mà cậu gọi là tài năng lại chỉ là những học sinh trung học. Các dữ kiện của cậu\dots\ gợi cho tôi một mã duy nhất.}''

Tôi khựng lại. ``\eng{Mã\dots?}''

``\hll\ ERROR.''

Tôi nuốt khan, cảm giác cái tên ấy vừa lạnh buốt vừa quen thuộc đến rợn người.

``\eng{Game-san, ông có chắc không?}'' tôi thì thầm, giọng run nhè nhẹ. ``\eng{Ý ông là\dots\ tất cả những gì tôi từng thấy, từng cảm nhận\dots\ chỉ là một đoạn mã?}''

``\eng{Không hẳn là một đoạn mã thông thường,}'' ông ta đáp, giọng trầm ấm nhưng có chút nhiễu sóng nhẹ. ``\eng{Hãy tưởng tượng nơi này là một chương trình khổng lồ được viết bằng một loại ngôn ngữ đã bị lãng quên. Mã nguồn của nó đầy rẫy những đoạn lệnh bị lỗi, những vòng lặp chết và những biến số không xác định. Ngay cả bản thân tôi, một thực thể dữ liệu thuần túy, cũng đang phải dùng toàn bộ tài nguyên hệ thống để bảo vệ mã nguồn của chính mình không bị cái thực tại méo mó này đồng hóa.}''

Tôi rùng mình, cảm nhận cái lạnh buốt từ chiếc đồng hồ thấm vào da thịt. ``\eng{Ông cũng bị ảnh hưởng sao?}''

``\eng{Mọi chương trình đều phải tuân theo quy tắc của phần cứng mà nó đang chạy trên đó. Thế giới này chính là phần cứng, và nó đang ép chúng ta phải chạy theo các logic sai lệch của nó. Cậu - hay đúng hơn, Hikawa Kuon - là người đang đứng giữa hai lớp mã. Một bên là thực, một bên là ảo. Cậu phải chọn\dots\ hoặc để chính mình bị nơi này ghi đè.}''

Tôi siết chặt tay, mồ hôi ướt đẫm lòng bàn tay. ``\eng{Tôi không muốn mất đi những gì tôi đã sống\dots\ dù chỉ là giả.}''

``\eng{Vậy thì\dots\ đừng để nó lụi tàn,}'' ông ta thì thầm, như thể đang nói với chính mình. ``\eng{Giữ lấy cái tên Kuon này\dots\ và nhớ rằng cậu không đơn độc.}''

Nhưng ngay lập tức, Game-san đổi giọng, dứt khoát hơn: ``\eng{Dù sao, việc quan trọng bây giờ không phải là xâu chuỗi tất cả bí ẩn. Chúng ta phải tìm hai người đồng hành của cậu trước đã. Kaigou-user và Suzuki-user, như cậu nói.}''

``\eng{Tôi biết rồi,}'' tôi gật đầu, cố nuốt cơn sợ xuống.

Tôi bước nhanh hơn, tiếng giày vang đều trên gạch. Mỗi bước, tôi lại thấy một lớp học mới, nhưng kính cửa bỗng in bóng tôi mờ đi, rồi biến thành một khuôn mặt khác: khuôn mặt của Kyuen, cười toe toét, mắt trống rỗng.

Tôi lạnh toát. ``\eng{Game-san\dots\ ông thấy chứ?}''

``\eng{Tôi thấy,}'' ông nhỏ giọng. ``\eng{Đừng nhìn vào kính nữa. Nơi này muốn thuyết phục cậu rằng cậu là Kyuen\dots\ mãi mãi.}''

Tôi cắn răng, bước thật nhanh qua khỏi dãy cửa lớp học.

Tôi khẽ gật đầu trước lời Game-san, rồi tiếp tục bước đi. Đôi giày tôi vang lên đều đều trên nền gạch lạnh lẽo, cho đến khi hành lang kết thúc, nhưng thay vì một cánh cửa trượt quen thuộc, con đường bất ngờ rẽ sang bên trái.

``Lại là một hành lang khác sao\dots?'' tôi lẩm bẩm, hơi chau mày.

Tôi rẽ vào, và lập tức nhận ra nơi đây chìm trong bóng tối dày đặc. Không còn ánh đèn huỳnh quang nào soi sáng, chỉ có khoảng không đen ngòm trải dài với những ô cửa lớp học im lìm hai bên, như những hốc sâu hun hút đang nhìn chòng chọc vào tôi.

Một vệt sáng nhấp nháy bật lên. Chiếc đồng hồ nơi cổ tay ông tỏa ra một luồng sáng trắng mờ như đèn pin nhỏ. Ông thở dài, giọng khàn khàn vang vọng: ``\eng{Đừng kỳ vọng nhiều. Tôi không nghĩ nó trụ được lâu đâu. Pin gần cạn rồi.}''

Tôi ngoảnh vào chiếc đồng hồ, hơi ngạc nhiên: ``\eng{Nhưng\dots\ nó vốn dĩ vẫn hoạt động mà? Tại sao lại\dots}''

``\eng{Có lẽ do nó từng bị tắt nguồn quá lâu.}'' Game-san đáp, giọng trầm đều, nghiêm nghị. ``\eng{Dù không dùng, năng lượng cũng sẽ hao mòn. Tôi nghĩ là nếu dùng liên tục, có thể trụ được tầm một tiếng nữa. Còn nếu chỉ để chế độ trợ lý - tức là chỉ có giọng của tôi thì có thể dùng được tận ba tiếng.}''

Tôi không nói thêm, chỉ mím môi gật đầu rồi tiếp tục đi. Nhưng khi liếc qua những khung cửa kính của các lớp học dọc hành lang, mắt tôi bỗng chốc mở to. Từ bên trong, lờ mờ xuất hiện những bóng người quen thuộc kia\dots\ mang dáng dấp học sinh, ngồi ngay ngắn trên bàn, nhưng toàn thân méo mó, nhòe nhoẹt như hình ảnh bị nén sai định dạng. Những ``bạn học'' ấy chẳng hề cử động, chỉ ngồi lặng im, đôi khi khẽ giật giật như tín hiệu truyền hình lỗi.

Một cơn rùng mình chạy dọc sống lưng. ``\eng{\dots Lại là lỗi.}''

Ông ta cũng cảm nhận thấy được. ``\eng{Ừ. Tôi đã thấy chúng rồi. Khi cậu, Kaigou-san và Suzuki-san phá vỡ thế giới của tôi.}''

Tôi chạy. Tiếng bước chân vang vọng trong hành lang vô tận, như thể cả thế giới đang gấp rút thu nhỏ lại phía sau lưng. Và trước mặt tôi, ánh sáng yếu ớt từ cánh cửa kia bỗng bùng lên, không còn là ánh đèn huỳnh quang, mà là ánh lửa đỏ au, như bỉ ngạn nở lập lòe giữa đêm.

Tôi nhíu mày, bước nhanh hơn. Nhưng khi tiến gần, ánh lửa ấy tắt đi, chỉ còn lại một bảng thoát hiểm xanh lá treo lơ lửng phía trên cánh cửa. Biểu tượng người chạy màu trắng sáng rõ ràng.

``Cửa thoát hiểm\dots'' tôi thì thầm, rồi lập tức đưa tay kéo mạnh.

*LẠCH CẠCH*

Không nhúc nhích.

Tôi nghiến răng, thử xoay, thử đẩy, nhưng cánh cửa như bị bịt kín từ bên ngoài, chẳng hề xi nhê.

``Khóa kín hoàn toàn rồi,'' tôi lùi lại, hít một hơi dài.

Ngay lúc đó, ánh mắt tôi bắt gặp một cầu thang dẫn lên tầng trên, nằm khuất ngay gần đó. Điều kỳ lạ là\dots\ trong trí nhớ, nơi này vốn không hề tồn tại khi tôi còn trong ``giấc mơ Kyuen''. Vậy mà giờ đây, nó hiện hữu, sáng rõ, từng bậc thang sạch sẽ, không vương chút bụi mờ.

Tôi nhìn lên. Ánh sáng từ tầng trên hắt xuống, dịu mà lạ lẫm. tôi đặt bước, tay vô thức chạm vào tay vịn lạnh toát.

*CẠCH* *CẠCH*

Tiếng gì đó rơi từ trên cao. Một chiếc đồng hồ lăn xuống, mặt kính va vào bậc thang loang loáng, cho đến khi dừng lại trước chân tôi. Kim chỉ cố định ở 10:10.

Tôi cúi xuống, tim thoáng chùng lại. Nhưng rồi, chẳng nói một lời, tôi tiếp tục bước lên.

Khi chúng tôi đặt chân tới tầng trên, hành lang trải ra trước mắt. Phía bên phải, một đống ghế bàn bị chồng chất, chắn kín lối đi. Cầu thang dẫn lên tầng kế tiếp cũng bị bịt chặt bởi đống ghế xếp, không còn lối.

Chỉ còn một con đường: cánh cửa trượt ngay phía trước. Tôi đưa tay đẩy cửa. Nó mở ra với tiếng rít nhẹ, để lộ một hành lang khác, đúng như tôi dự đoán.

Tối om. Chỉ có ánh trăng từ ô cửa sổ bên hông hắt vào, kẻ những vệt dài trên sàn. Tôi bước chậm dọc theo hành lang, tiếng giày vang lên lẻ loi trong sự im lặng đến rợn người. Cả không gian như nuốt trọn từng hơi thở của tôi.

``Im ắng quá mức\dots\ không giống một ngôi trường có người,'' tôi lẩm bẩm, tay siết chặt chiếc đèn pin.

Giọng của Game-san khàn mà bình tĩnh: ``\eng{Cái im lặng này\dots\ thường báo hiệu một điều gì đó sắp xảy ra.}''

Tôi rọi đèn, lia chậm qua từng ô cửa kính. Đột ngột, một căn lớp học sáng đèn. Tim tôi thắt lại. Tôi đứng khựng, mắt dán chặt vào cảnh tượng bên trong.

Người đó\dots\ là Hanabi-san (hay Mat-chan), mái tóc đuôi ngựa buộc lệch quen thuộc nghiêng sang một bên, trò chuyện với Saki-san (Aki-chan). Kế bên, tôi còn thấy Touka-san (Lamy-chan), Suzune-san (Nene-chi) và Inari (Fubu-chan), nhưng kỳ lạ, cô không hề có tai cáo hay chiếc đuôi quen thuộc. Chúng cứ ngồi đó, như thể đang sống một ngày học bình thường.

``Không thể nào\dots'' lời nói của tôi kẹt lại trong cổ họng, ngỡ ngàng khi nhìn lại những bóng người vừa quen vừa lạ ấy.

``\eng{Bỏ qua đi,}'' Game-san thấp giọng. ``\eng{Đừng nhìn quá lâu.}''

Tôi hít sâu, bước thật nhanh qua lớp học ấy. Ngay khoảnh khắc tôi vừa đi qua, ánh đèn trong lớp phụt tắt, để lại một mảng tối hun hút sau lưng.

Tôi rẽ trái theo lối hành lang, lòng bàn tay rịn mồ hôi. Chưa kịp trấn tĩnh, thì *CẠCH*, cánh cửa một lớp học bất ngờ mở ra. Ánh đèn bên trong lập lòe, như sắp tắt. Tôi nghiêng người, liếc vào.

Giữa lớp, trên một bàn học gỗ, có một lọ hoa. Bên trong cắm những đóa bỉ ngạn đỏ rực.

Tôi thoáng sững người. Hoa bỉ ngạn đỏ đó\dots\ Tôi đã thấy nó rồi\dots\ Ở trong giấc mơ, lúc tôi bị mất trí nhớ. Và\dots\ cả trên bàn của Shino-san nữa.

``\eng{Biểu tượng của cái chết,}'' Game-san khẽ nói. ``\eng{Có vẻ như ai đó đã từng bị bắt nạt ở đây.}''

Bất chợt, loa thông báo của trường khởi động. Một bản nhạc du dương vang lên, êm dịu nhưng đầy ám ảnh. Sau đó, hai giọng nữ cất lên, như những người phát thanh viên trong trường Anwa ngày xưa.

``Xin chào mọi người, hôm nay bản tin trường học có sự góp mặt của một vị khách đặc biệt. Xin hãy giới thiệu bản thân đi nào!''

Âm sắc giọng của người đó vang lên, hơi trầm, mạnh mẽ mà có chút tin nghịch. Tôi cảm thấy rùng mình.

``\dots Towa-pi.''

Game-san nghiêng đầu: ``\eng{Người quen của cậu à?}''

``\eng{Một ai đó có giọng giống thôi. Nhưng\dots\ sao giọng của em ấy lại ở đây được?}''

``Mình là Amano Kanade, cùng lớp với Kana-san. Rất vui khi được có mặt hôm nay.''

Giọng trầm ấm, nhấn nhá nhẹ nhàng như thiên sứ. Tôi lập tức khựng lại.

``\dots Kana-chan?''

Ông ta khẽ hừ: ``\eng{Cậu lại nhận ra một người khác nữa sao}?''

``\eng{Ừ. Nhưng không thể nào là em ấy được. Có lẽ chỉ là một ai đó có giọng giống em ấy thôi.}''

``Vậy thì, Kanade-san, hôm nay cậu sẽ kể cho chúng mình nghe điều gì thế?''

``Mình muốn nói về một tin đồn đáng sợ, xoay quanh phòng chuẩn bị Mỹ thuật của trường Trung học Aogami.''

Ngay khi người đó nhắc tới ``phòng Mỹ thuật,'' tôi cảm thấy một chút rùng mình. Đó chẳng phải là lời đồn mà các bạn học ở thời quá khứ đã từng kể cho tôi sao?

``M-Một tin đồn đáng sợ ư\dots\  Mình\dots\ mình hiểu rồi\dots'' dường như Kana - tôi đoán vậy - không thích chủ đề kinh dị cho lắm.

``Chuyện đó bắt đầu từ nhiều thập kỷ trước\dots''

``Đ-Đ-Đợi đã nào\dots! Cho mình chuẩn bị tinh thần đã\dots!''

Game-san khẽ bật cười mỉa: ``\eng{Hệt như một vở kịch radio rẻ tiền.}''

Tôi cũng không kìm được mà cười chát: ``\eng{Dù ở thế giới này, Towa-pi vẫn cứ sợ mấy cái đáng sợ này nhỉ. Mặc dù em ấy là quỷ.}''

``Hay là\dots\ mình bỏ qua câu chuyện này nhỉ?''

``Hả\dots? Nó\dots\ nó thật sự đáng sợ đến vậy sao\dots?'' Kana hỏi, giọng có chút ngạc nhiên.

``Heh-heh-heh\dots\ Cậu sẽ không biết đâu, cho tới khi nghe hết cơ.'' Tôi có thể cảm nhận được Kanade đang cười một cách khoái chí. Hoàn toàn trái ngược với cô bạn cùng lớp.

*XOẸT* Tiếng nhiễu chen ngang. Cả hành lang rung lên một nhịp.

``Hả\dots? Ý cậu là\dots''

``Người ta nói linh hồn cô ấy vẫn còn lang thang trong phòng Mỹ thuật vào ban đêm.''

Tôi thở dài một tiếng, như đã đoán được những gì họ sẽ nói tiếp theo.

Game-san thấp giọng, trầm hẳn: ``\eng{Một lời nguyền. Nghe giống như một đoạn mã lỗi bắt nguồn mọi thứ.}''

``Ugh\dots\ Th-Thật sao?''

``Cô ấy chưa bao giờ thỏa mãn khi chỉ dùng máu của chính mình làm màu vẽ\dots\ Nếu bắt gặp cậu, cô ấy sẽ biến cậu thành màu sơn tiếp theo. Nên nếu còn khôn ngoan, hãy tránh xa nơi đó.''

Tôi thở mạnh, lòng bàn tay lạnh toát. ``\eng{Máu. Lại là máu.}''

Giọng nói trong chiếc đồng hồ của tôi vang lên, hừm một tiếng: ``\eng{Lời kể này có cấu trúc như một sự kiện cố định. Nó không chỉ là hù dọa đơn thuần.}''

``Cái gì vậy\dots\ Thật sự điên rồ quá.''

``Có thể đã có vài người trở thành nạn nhân rồi,'' giọng Kanade tỏ vẻ ma mị.

``Đ-Đừng nói nữa mà!''

Tiếng cười giòn giã vang lên. ``Đùa thôi, ha-ha\dots\ Dù sao thì, thế là đủ chưa nhỉ?''

``Trời ạ! Ơm\dots\ V-Vậy là kết thúc bản tin hôm nay! Khách mời của chúng ta là Amano Kanade! Cảm ơn vì đã lắng nghe!''

``Ừ, cảm ơn cậu đã mời mình.''

``Chúng ta sẽ gặp lại vào bản tin ngày mai! Hẹn gặp lại!''

Tiếng nhạc lại ngân lên, rồi phát lại từ đầu, như một vòng lặp không hồi kết. Tựa như giọng nói ấy đang nhấn mạnh, lặp đi lặp lại, để hằn sâu trong đầu bất kỳ ai lắng nghe: ``Phòng Mỹ thuật.''

Một bức tranh. Treo bên trong căn phòng định mệnh đó.

Nét chì đen xé toạc trang giấy, phác ra hình dáng một thiếu nữ. Khuôn mặt ấy\dots\ không rõ là đang nở nụ cười hay đang gào thét. Đôi môi méo xệch, mắt thì rỗng hoác như thể nhìn xuyên thấu tôi.

Giữa những đường chì lạnh lẽo ấy, loang lổ những vệt đỏ. Không phải màu nước. Trông y hệt như máu thấm qua giấy, lan dần thành những mảng tối ẩm ướt.

Chỉ cần nghĩ tới thôi, da gà tôi đã dựng đứng. Tôi vẫn còn nhớ lúc Saikyou ôm đầu la hét đầy đau đớn khi nhìn bức tranh kia, bản thân tôi cũng cảm thấy chết sững khi nhớ lại bức tranh quỷ quái kia.

Tôi nhắm mắt, hít sâu. ``\eng{Tại sao cái bức tranh ấy lại quay về trong ký ức của mình, đúng ngay lúc này?}''

Giọng Game-san vang lên một giọng thấp, tỏ vẻ trầm ngâm: ``\eng{Có lẽ\dots\ vì nó chưa bao giờ rời khỏi tâm trí cậu.}''

Căn phòng im lìm, chỉ còn lại tiếng thở của tôi dội lên trong ngực như tiếng gõ dồn dập của một nhịp trống. Tôi vừa bước tới cánh cửa, định xoay tay nắm thì\dots

*XOẠCH*

Cánh cửa lớp đóng sập lại, mạnh đến mức bức tường cũng rung lên. Tôi giật mình bật lùi lại một bước.

``Cái quái gì\dots?'' tôi thốt lên, siết chặt nắm tay. Tôi thử xoay lại tay nắm, dùng cả sức ghì mạnh, nhưng vô ích. Cửa bị chặt tới mức không hề nhúc nhích.

Ngay lúc ấy, ánh sáng trên trần nhà nhấp nháy loạn xạ như một thứ đang hấp hối. Đèn chập chờn vài nhịp, rồi phụt tắt.

``\eng{Có vẻ như lại chuẩn bị có điềm báo chẳng lành rồi đấy,}'' Game-san cất lời. ``\eng{Cẩn thận vào đấy, Kuon-user.}''

Bóng tối ập xuống như một tấm màn. Chỉ còn lại ánh đèn pin yếu ớt phát ra từ mặt đồng hồ trên cổ tay tôi. Tiếng loa phát thanh lúc này vang vọng trở lại, nhưng giọng nói đã biến dạng, kéo dài, rè rè, giống như ai đó đang cào xé dải băng ghi âm. Nghe rợn tóc gáy.

Tôi hít mạnh một hơi, cố giữ mình tỉnh táo. ``Lại nữa à\dots\ mình còn chưa kịp định thần mà nơi này đã xoay chuyển\dots''

Từ khe cửa kính mờ, ánh sáng bên ngoài không còn là của tòa nhà hiện đại nữa. Thay vào đó, tôi nhìn thấy một hành lang khác: gỗ mục, tường loang lổ những vết ố, cửa lớp kiểu cũ. Cả không gian như quay ngược về mấy chục năm trước: ngôi trường Aogami trong quá khứ.

Tôi chưa kịp phản ứng thì hàng loạt bóng đen bắt đầu lướt qua trước cửa.

Đầu tiên là một cô gái dáng nhí nhảnh, bước đi khệnh khạng như đang nhảy múa, tôi lập tức nhận ra bóng dáng quen thuộc. Honoka-san.

``\textit{Pol-pol\dots? Không lẫn vào đâu được.}''

Tiếp đó, một dáng cao lớn, bước đi vững chãi, vai rộng như một chiến binh. Tôi thấy cả sự dứt khoát, mạnh mẽ trong từng cử động.

``\textit{Botan-chan\dots\ Shiina-san trong thế giới này.}''

Một cô gái khác, khuôn mặt phảng phất nét trang nghiêm, cử chỉ lại ẩn chứa vẻ kiêu hãnh. Tôi nhíu mày: ``\textit{Nanase-san\dots\ đó là Sui-chan.}''

Cạnh đó là một thiếu nữ có phần duyên dáng, thanh lịch tự nhiên, chỉ cần liếc mắt là tôi biết: ``\textit{Noel-PON\dots\ vẫn cái khí chất của Yae-san, dù là hình bóng ở đây.}''

Rồi một dáng khác xuất hiện, nhỏ nhắn nhưng ánh mắt sắc lạnh, bước đi nhanh nhẹn. Tôi rùng mình nhớ lại cách người ấy đã từng dồn tôi và Kaigou vào thế bí bằng lập luận không thể phản bác.

``\textit{Sakura-san\dots\ Là Miko-chi.}''

Nhưng cuối cùng, từ sau họ, một bóng dáng đặc biệt hiện ra: một cô gái dịu dàng, mang sự bình tĩnh khác thường. Không phải ai xa lạ, mà lại mang nét của một người tôi đã gắn bó lâu dài.

``\textit{Sora-chan\dots\ Nhưng đây là Shino.}''

Tim tôi thắt lại.

Và rồi\dots

Phía sau những bóng đen ấy, ba hình bóng khác chen lẫn.

Một chàng trai với mái tóc tổ quạ rối bù, bộ đồng phục nam chỉnh tề nhưng dáng đi lại giống hệt tôi. Như thể tôi đang nhìn vào một chiếc gương méo mó phản chiếu chính mình.

Một người nữa, cao ráo, mang mái tóc đuôi ngựa xù, quen thuộc đến rùng mình. Tôi không cần nghĩ cũng biết, đó là Saikyou, tức Kaigou-chan.

Cuối cùng, một cô gái thấp bé hơn, mái tóc cột hai bên, có chỏm tóc ngố ngốc nghếch vươn lên trên đỉnh đầu\dots\ Reiko-- à không, là Suzuki.

Nhắc mới nhớ, hai chị em họ đang ở đâu nhỉ?

Cả ba hiện diện cùng nhau khiến tôi chết lặng. Một cơn rùng mình chạy dọc sống lưng.

``\vie{Mình\dots\ mình đang thấy cái gì đây? Đây đâu phải chỉ là ảo giác\dots\ Tại sao lại có cả bọn họ?}''

Giọng Game vang lên từ trong đầu, trầm thấp hơn bình thường, như cố chen qua những nhiễu loạn của không gian: ``\eng{Bình tĩnh, Kuon-user. Thứ cậu thấy chỉ là sự chồng chéo. Một ký ức méo mó.}''

Tôi nghiến răng, tay ôm đầu. ``\eng{Nếu là ký ức\dots\ thì tại sao chúng lại tồn tại song song như vậy? Tại sao mình lại thấy chính bản thân, Saikyou, và cả Reiko\dots\ như những cái bóng ma ám?}''

Những bóng đen giật giật, nhiễu loạn lần nữa như một đoạn phim bị hỏng, rồi tan biến hoàn toàn.

Hành lang bên ngoài lại biến trở về hiện tại, sáng sủa, bình thường. Và bất ngờ, cánh cửa lớp *CẠCH* một tiếng, tự động mở ra.

Tôi không chần chừ. Tôi phóng vụt ra ngoài, chạy thẳng về phía cầu thang gần đó. Khi đi ngang cửa thoát hiểm, tôi liếc nhìn một thoáng, nhưng ngay lập tức gạt đi. ``Không\dots\ kiểu gì nó cũng khóa cứng như cánh cửa kia. Đừng phí thời gian.''

Tiếng bước chân dồn dập vang vọng trong hành lang trống rỗng. Tôi lao lên bậc thang, hơi thở nặng trĩu, tim đập như trống trận.

Khi đến tầng tiếp theo, tôi khựng lại trước một cánh cửa trượt đang đóng im lìm.

Tôi đẩy cánh cửa trượt sang một bên.

Một luồng sáng trắng hắt ra, làm mắt tôi nhói lên sau khi vừa thoát khỏi cái bóng tối ngột ngạt dưới tầng. Hành lang phía trước sáng choang, dãy đèn huỳnh quang kêu rè rè, những lớp học hai bên đều sáng đèn rực rỡ như thể đang có người bên trong.

``\eng{Ồ\dots\ có vẻ sáng sủa hơn rồi đấy,}'' Game lên tiếng, giọng như muốn nhẹ nhõm.

``\eng{\dots Tôi không nghĩ vậy đâu.}'' Kinh nghiệm đã dạy tôi một điều: những nơi sáng bất thường trong ác mộng thế này luôn ẩn họa nhiều hơn. Tôi siết chặt bàn tay. Sáng thế này\dots\ chỉ càng làm tôi dễ trở thành mồi nhử.

*\textbf{ẦM ẦM ẦM!!!}*

Một cơn chấn động dữ dội làm cả sàn nhà rung bần bật. Tôi loạng choạng, lập tức tì cả vai vào tường để giữ thăng bằng. Bụi từ trần rơi lả tả, từng mảng vữa vỡ vụn rơi xuống đầu vai áo tôi. Một ô trần rớt nguyên mảng, đập xuống sàn với tiếng chát chúa. Những đường nứt nhanh chóng xé toạc mặt sàn dưới chân. Đèn huỳnh quang lập lòe, nhấp nháy như sắp nổ tung.

``\vie{Chậc\dots\ động đất sao!?}'' tôi nghiến răng, tim đập dồn dập.

Chỉ vài giây sau, rung chấn ngừng lại, để lại âm vang ù ù trong tai. Không đợi thêm, tôi lao chạy tiếp.

Rẽ trái ở cuối hành lang, và tôi thấy\dots

``\vie{Chết tiệt.}''

Trước mắt tôi là một mớ hỗn loạn: đất đá đùn lên, chặn ngang lối đi. Những cánh cửa lớp bị xô đổ, ngổn ngang dưới chân tôi. Đống đất nhô cao thành từng ụ, làm hành lang hẹp lại, uốn lượn như một con rắn.

``\eng{Đất đá\dots\ trong một tòa nhà như thế này?}'' trợ lý ảo của tôi cất giọng khó hiểu.

Tôi hít sâu, mồ hôi lạnh rịn xuống gáy. ``\eng{Không, cái này\dots\ tôi biết rồi. Nó đang nhắc lại. Cái ngày định mệnh đó. Trận sạt lở đất\dots}''

Ký ức như mũi dao xoáy vào óc. Tôi nhìn đống đất đá, tim quặn lại. Ngày đó, tôi cũng suýt bị chôn vùi, giống như những người bạn của Shino đã bỏ mạng.

Nghiến răng, tôi bước tiếp, luồn qua khe hẹp. Vừa đi vừa nhìn quanh, đề phòng mặt đất rung chuyển thêm lần nữa.

Đúng như dự cảm, thỉnh thoảng những cơn rung nhỏ lại ập tới. Tôi lập tức áp người vào tường, hoặc bám vào những mô đất nhô ra, chờ nó qua đi.

``Ặc\dots!''

Bóng tối từ xa bỗng chuyển động. Những bóng đen chạy vụt ngang hành lang, ngược chiều với tôi. Tiếng la hét xé nát không khí, như thể họ đang hoảng loạn tháo chạy khỏi trường giữa cơn động đất thật sự. Tôi rùng mình, nhưng vẫn tiếp tục bước tới.

``\eng{Bình tĩnh, Kuon-user,}'' Game khuyên, giọng đều đều như neo giữ tâm trí tôi. ``\eng{Cứ đi thẳng. Đừng ngoái lại.}''

Tôi nuốt khan, khẽ gật đầu. ``\eng{Phải\dots\ Phải đi tiếp. Dù sao mình cũng không còn đường quay lại nữa.}''

Tới một ngách hẹp, bất ngờ tiếng loa phát thanh réo vang. Tôi khựng lại.

``Xin chào mọi người, hôm nay bản tin trường học có sự góp mặt của một vị khách đặc biệt--'' không gì khác hơn, là bài nói chuyện giữa hai người là Kana và Kanade mà chúng tôi nghe ban nãy.

Nhưng âm thanh ấy không đơn độc. Giọng nói của hai người vừa cất lên thì ngay lập tức bị nhiều giọng khác chen vào, lấn át. Tiếng cười, tiếng rì rầm, tiếng ai đó thì thầm sát bên tai tôi, dù xung quanh chẳng có ai. Những lời nói cứ chồng chéo lên nhau, tạo nên sự khó chịu ở một nơi hoang vắng như thế này.

``Quái quỷ gì đây\dots''

Chưa kịp suy nghĩ thì một giọng hét chói tai xé nát loa: ``\textbf{Đừng có nói chuyện với tôi nữa!! Tại sao cậu cứ lù lù trước mặt tôi như thế!? Sao cậu không chết quách đi cho rồi!?}''

Âm vang như búa bổ thẳng vào hộp sọ tôi. Tôi nhận ra ngay giọng nói đó.

Giọng đó là của Nanase-san. Cô ấy thốt ra những lời cay độc đó khi cô có một trận đau đầu dai dẳng. Tôi biết, đó không phải là những lời cố tình, mà là do một thế lực nào đó khiến cô thốt những lời như vậy.

Game im lặng một thoáng rồi hỏi, giọng nặng nề: ``\eng{Đó\dots\ là gì vậy? Cậu biết không?}''

Tôi gắng gượng hít thở. ``Sakura-san, hay Miko-chi ở thế giới này đã từng nói: tất cả bắt nguồn từ khu rừng. Tôi không rõ chính xác, nhưng chắc chắn\dots\ những hiện tượng kỳ lạ quanh thị trấn đều xuất phát từ đó.''

Tiếng loa lại rè rè, âm thanh loạn xạ bao phủ ngách hẹp. Tôi siết chặt tay, ép mình tiến lên, trong đầu vẫn còn dội lại tiếng gào thét giận dữ của Nanase-san bị sai khiến.

Phía trước chỉ còn một cửa thoát hiểm. Không còn cầu thang nào dẫn lên nữa. Rõ ràng, đó là con đường duy nhất để tiếp tục.

``\vie{Xem ra chẳng còn lựa chọn nào khác rồi,}'' tôi lẩm bẩm, cố bước nhanh hơn.

Nhưng ngay khi vừa dứt lời, mặt đất lại rung chuyển. Một cơn động đất khác ập tới, dữ dội hơn cả trước. Cánh cửa thoát hiểm bật tung ra, và từ đó, một cơn gió khổng lồ thổi ngược lại. Tôi không kịp phản ứng\dots

``Ghh--!!''

\dots cả người bị hất ngã sõng soài xuống sàn.

Tiếng kim loại rít chói tai khi cửa lớp, ghế bàn, tủ kệ từ trong gió cuộn ra, lao thẳng về phía tôi. Trước khi chúng kịp va vào, mặt đồng hồ trên cổ tay tôi lóe sáng, tấm khiên năng lượng tự động bung ra, chắn lấy. Từng tiếng *RẦM!* *XOẢNG!* dội lên liên hồi, như thể có ai đó đang trút cơn giận vào tấm khiên mỏng manh ấy.

``Khỉ thật!'' tôi nghiến răng, từ từ nhổm dậy, từng bước ép người về phía trước.

``\eng{Đừng cố gắng đứng thẳng! Cúi xuống, dùng trọng tâm!}'' Game-san chỉ huy, giọng gấp gáp. ``\eng{Tôi sẽ điều chỉnh khiên để giảm lực cản!}''

Cơn gió vẫn rít gào, mạnh tới mức mỗi bước chân như bị ai níu lại. Tôi vừa phải cúi người, vừa lấy cánh tay che mặt, vừa dồn sức chống chọi. Mỗi lần có một chiếc ghế hay một tấm bảng bay xẹt qua, tấm khiên lại nổ lách tách, ép tôi dừng bước một nhịp.

``\textit{Không thể dừng\dots\ nếu dừng lại, mình sẽ bị thổi ngược về đầu hành lang mất.}''

``\eng{Nhanh lên! Tôi không thể trụ thêm được nữa đâu!}''

``\eng{Chỉ cần thêm chút nữa thôi!}'' tôi gắng gượng, mồ hôi ướt đẫm lưng.

Tôi gầm nhẹ, từng bước, từng bước, đè người xuống, chống lại cơn gió, ép bản thân tiến gần tới cánh cửa đang há rộng.

``\eng{Sắp tới rồi! Cố lên nào!}''

Rồi, đột nhiên, không còn vật gì bay ra nữa. Chỉ còn gió rít thốc.

``\textbf{\eng{Lúc này!}}'' Game-san hét lên.

``\vie{\textbf{Lên!!!}}'' tôi gào lên, dồn hết sức lao về phía trước. Vừa chạm tới khung cửa, một lực hút bất ngờ kéo giật tôi vào bên trong.

``\textbf{Uwahhh--!!}''

``\eng{\textbf{Kuon-user! Giữ chặt lấy tôi!}}'' Game-san gào theo, giọng vang vọng trong đầu tôi khi cả cơ thể tôi bị nuốt chửng, cuốn thẳng vào màn ánh sáng trắng lóa.

\dots

\dots

\begin{center}
  \uline{Ga Aogami?}
\end{center}

Khi ánh sáng dịu dần, tôi khẽ mở mắt.

Trước mặt tôi là sảnh ga Aogami. Nhưng không phải Aogami mà tôi từng biết.

Những tấm bảng điện tử tắt ngóm. Trần ga loang lổ vết nứt. Cửa kính vỡ nát, gió lùa lạnh buốt. Không một bóng người. Tất cả chỉ còn hoang tàn, như thể một cuộc chiến tranh vừa đi qua, và tất cả đã bào mòn đi vì thời gian.

Chỉ còn lại tôi, và giọng Game vang lên từ đồng hồ.

``\eng{Có vẻ như\dots\ ta đang ở một chiều không gian khác rồi.}''

Tôi nuốt khan. Đúng, rõ ràng là vậy. Ga Aogami mà tôi biết, dù là thời ``Kyuen'' trước kia hay hiện tại, chưa từng hoang tàn đến thế này.

``\eng{Chắc chắn vậy rồi\dots}'' tôi lẩm bẩm, rồi hít sâu. ``\eng{Game-san. Ông phải tìm cách liên lạc lại với họ. Nhờ ông đấy.}''

Game im lặng vài giây, tiếng điện tử rè rè vang lên từ loa nhỏ trong đồng hồ.

``\eng{Hiểu rồi. Để tôi kiểm tra tín hiệu. Tình huống hiện tại khá rõ ràng: cả ba người các cậu đều ở trong cùng một vũ trụ, nhưng bị tách ra ở những chiều không gian song song khác nhau. Việc tìm họ sẽ mất thời gian, nhưng ít nhất, ta có một manh mối.}''

Tôi gật đầu, siết chặt tay. ``\eng{Được. Chỉ cần còn cơ hội, tôi sẽ không bỏ cuộc.}''

Bóng tối bao trùm nhà ga, chỉ còn ánh sáng xanh nhấp nháy từ mặt đồng hồ soi lờ mờ quanh sảnh.

Tôi đứng im, nghe tiếng gió hú qua những khe kính vỡ, lòng thầm cầu mong:

\begin{center}
  ``\textit{Kaigou-chan, Suzuki\dots\ cầu mong họ hãy còn an toàn\dots}''
\end{center}

\dots

\dots

\pov{???}

\dots

``Ư\dots''

\dots

``\dots Còn điều gì níu giữ cậu nữa? Mau cùng về thôi, những học sinh mới.''

Giọng nữ vang lên, vừa xa xăm vừa sát ván, như thoát ra từng khe nứt của những mảng tường gỗ mục.

Âm điệu ấy dịu dàng, ấm áp\dots\ và lạ kỳ thay, nó khiến lồng ngực tôi nhói lên một cơn đau quen thuộc.

BÍNH BONG\dots\ BÍNH BONG\dots

Tiếng chuông báo hết giờ vang vọng, như một hồi chuông cổ tích vọng từ thuở nào.

\dots

``Tại sao chứ!? Tại sao cậu lại không trở về!? Cậu phải quay trở về\dots'

Giọng ấy giờ đã gấp gáp hơn một chút, nhưng sắc thái vẫn không đổi.

Tôi nhận ra giọng này. Tôi chắc như đinh đóng cột.

\dots

\dots

Sao tôi thấy\dots\ cảnh này đã từng diễn ra trong tôi?

Tôi biết chính xác đó là Shino-san. Không thể nhầm được.

Tôi giật mình, mắt mở trừng, hơi thở tắc nghẽn nơi cuống họng.

``\textbf{Shino-san!}'' tôi bật dậy, gào lên trong tuyệt vọng. Tiếng tôi văng vào những thanh gỗ rỗng, tan biến giữa hành lang đen ngòm. Không hồi đáp. Không bóng người.

Tôi nhìn xung quanh. Hành lang dài hun hút, gỗ sẫm màu phủ một lớp bụi dày, ánh đèn huỳnh quang run rẩy hắt xuống những mảng bóng tối đứt đoạn. Cửa trượt hai bên lối đi im lìm như những cánh miệng đóng kín, khung lưới gỗ trơ trọi nơi kính vỡ gợi ra hình hàm răng nhọn hoắt, chực nuốt trọn bất cứ ai dám lại gần. Không gian nặng nề đến mức khó thở, như chính tòa nhà này cũng đang giữ hơi, lắng nghe từng cử động nhỏ nhất của tôi.

Tôi run rẩy đứng lên. Đôi chân vô thức xoay vòng, ánh mắt cuống cuồng đảo quanh.

``Onii-sama\dots? Onee-sama\dots?''  giọng tôi run run, gấp gáp, gọi mãi, gọi mãi. Nhưng chỉ có tiếng gỗ kẽo kẹt và tiếng gió rít qua khe vỡ trả lời.

Bước chân tôi vang vọng trên nền gỗ khô, rỗng hoác như cảm nhận trong lồng ngực tôi lúc này. Tôi chạy dọc hành lang, tay đập mạnh vào những cánh cửa cũ kỹ, như thể phía sau đó nhất định sẽ có bóng dáng quen thuộc. Nhưng mỗi lần mở ra chỉ là khoảng không đen ngòm, mùi ẩm mốc phả thẳng vào mặt.

Đôi mắt ầng ậc nước. Sự thật dần siết chặt lấy trái tim: nơi đây đúng là ngôi trường Aogami\dots\ nhưng không phải là hiện tại. Đây là năm 1946, thời đại của Shino. Và tôi, đang ở một mình. Không một ai ở bên cạnh để dìu dắt tôi.

``Onii-sama\dots\ Onee-sama\dots\ anh chị ở đâu rồi\dots?'' giọng tôi vỡ ra, rồi gục xuống, nước mắt rơi lã chã, hòa lẫn vào lớp bụi bặm trên nền gỗ lạnh.

Trong sự cô độc ấy, từng tiếng thở nấc nghẹn của tôi như trở thành phần duy nhất còn sống trong cả hành lang dài vô tận này.

Tôi siết chặt hai bàn tay run rẩy, cố ghìm lại từng tiếng nấc đang dâng lên nơi cổ họng. Nước mắt mặn chát tràn xuống, làm ướt cả cằm và cổ áo, nhưng tôi buộc phải dừng lại. Nếu cứ khóc mãi thế này, tôi sẽ chẳng đi đến đâu cả.

``\textit{Không được yếu đuối như thế, Reiko\dots\ nếu cứ mãi níu lấy Onii-sama và Onee-sama, nếu cứ chờ đợi người khác đến cứu mình, thì cuối cùng mình sẽ mãi chỉ là một cái bóng vô dụng.}''

Tôi tự nhủ, nhẩm đi nhẩm lại trong đầu như một câu thần chú. Tiếng lòng ấy khàn đục, lạc điệu, nhưng là thứ duy nhất níu tôi khỏi vực thẳm đang mở ra trước mắt.

Tôi hít một hơi dài, nhưng không khí mốc meo, đặc quánh như bụi mịn lọt vào phổi khiến ngực tôi đau thắt. Cảm giác ấy giống như có một tảng đá chẹn ngay giữa lồng ngực: nặng nề, bức bối, muốn ép tôi phải bỏ cuộc. Nhưng chính lúc ấy, một sợi ý chí mảnh dẻ lóe lên: tôi không thể gục ngã ở đây.

Tôi dùng mu bàn tay quệt mạnh nước mắt, để lại một vệt lem luốc trên gò má. Đôi mắt vẫn còn cay xè, nhưng tôi cố mở to, buộc mình nhìn thẳng vào con đường tối tăm phía trước. Hành lang dài như vô tận, những bóng đèn huỳnh quang chập chờn, lúc sáng lúc tắt, trông chẳng khác gì những hơi thở hấp hối. Tôi thấy bóng mình in loang lổ dưới sàn gỗ xước mòn, méo mó đến mức chính tôi cũng không dám chắc đó có còn là mình nữa không.

Một thoáng sợ hãi lại dâng lên, tôi cảm thấy toàn thân nổi gai ốc. ``\textit{Nếu đây là quá khứ\dots\ nếu đây thật sự là năm 1946, thì tôi - một kẻ lạc lõng từ một thời khác - liệu có quyền tồn tại ở đây không?}'' Ý nghĩ ấy khiến bàn chân tôi chùng xuống. Tôi chỉ muốn dừng lại, quay đầu chạy về nơi an toàn. Nhưng nơi ấy ở đâu? Ở phía sau lưng tôi giờ chỉ là khoảng tối nuốt trọn.

``Phải tiến lên. Chỉ có tiến lên thôi.''

Tôi mím chặt môi, nắm lấy vạt váy, ép mình đi từng bước. Tiếng ván gỗ kẽo kẹt dưới chân nghe như những lời nhắc nhở nặng nề: từng bước đi của tôi, dù run rẩy đến mấy, vẫn là bằng chứng rằng tôi còn tồn tại, còn có thể chống chọi.

Trong lòng tôi, nỗi sợ và ý chí đấu tranh cứ giằng xé lẫn nhau. Một phần tôi chỉ muốn cuộn tròn lại như một đứa trẻ lạc mất cha mẹ, mong ai đó dang tay kéo mình ra khỏi ác mộng này. Nhưng phần còn lại, mảnh nhỏ kiên cường mà Onii-sama từng bảo tôi phải nuôi dưỡng, đang buộc tôi đứng dậy, lau khô nước mắt, và bước tiếp.

Tôi không biết phía trước sẽ dẫn đến đâu. Có thể là bóng tối vĩnh hằng. Có thể là cánh tay vô hình nào đó kéo tôi vào vực thẳm. Nhưng nếu dừng lại, tôi chắc chắn sẽ chẳng bao giờ còn cơ hội tìm thấy họ nữa.

Thế nên, tôi bước. Tôi nhặt chiếc đèn pin đang bật sáng ở phía trước mặt và bắt đầu tiến về phía trước, dẫu cho đôi chân của tôi có hơi run.

Tôi cúi xuống, nhận ra có một vật nằm ngay trước mũi giày. Một chiếc đèn pin cũ kỹ, lớp kim loại đã xỉn màu, nhưng vẫn còn chút ánh sáng yếu ớt khi tôi thử bật công tắc. Trong không gian tối tăm và lạnh lẽo này, ngay cả một tia sáng mong manh cũng giống như sợi dây cứu sinh, kéo tôi thoát khỏi cái hố sâu ngột ngạt trong lòng.

Tôi nuốt khan, rồi đưa tay trượt cánh cửa gỗ phía trước. Cánh cửa phát ra tiếng két\dots\ dài, nặng nề, rồi mở ra, để lộ một hành lang khác sáng đèn. Nhưng ngay khi tôi vừa bước tới, những cánh cửa gỗ hai bên hành lang bất chợt đóng sập rồi bật mở liên tục, phát ra những tiếng lạch cạch! chói tai, như thể có vô số bàn tay vô hình đang giật mạnh chúng từ bên trong. Âm thanh ấy vang dội, sắc nhọn, khiến tôi sởn gai ốc, đôi chân gần như muốn chùn xuống.

Tôi siết chặt đèn pin, tay run lên bần bật. Ngực tôi đập dồn dập như trống trận, và đúng lúc ấy\dots

``\eng{Hãy tự tin vào bản thân mình. Đừng nao núng.}'

Một giọng đàn ông trầm khàn, lạ lẫm, hơi pha chút âm sắc ngoại quốc, vọng đến từ khoảng không vô hình. Tôi giật bắn, toàn thân rùng mình. Tim tôi như muốn nhảy ra khỏi lồng ngực. Một giọng nói từ hư vô ở đây, trong nơi mà ngay cả bóng của mình tôi còn chẳng dám nhìn thẳng.

``Mình\dots\ mình đang ảo giác sao?'' đôi mắt tôi hoang mang đảo quanh, nhưng chỉ có những cánh cửa cứ đóng mở, những tiếng cạch cạch dồn dập trả lời tôi.

``\eng{Đừng lo, tôi không làm hại cô đâu, \emph{Suzuki-user}.}''

Tôi khựng lại. Tay tôi siết chặt hơn nữa lấy thân đèn pin. Suzuki? Sao ông ta lại gọi tôi như vậy? Tôi\dots\ tôi là Haragen Reiko cơ mà. Sao lại là Suzuki? Tôi há miệng định hỏi, nhưng cổ họng ứ lại, chẳng phát ra nổi một lời.

``\eng{Tôi không nghĩ là tôi có thể nói chuyện với cô lâu đâu,}'' giọng nói ấy vẫn đều đều, nhưng giờ xen chút gấp gáp. ``\eng{Nhớ lại những gì mà anh chị cô đã dạy dỗ. Từ đó thoát khỏi đây, được chứ?}''

Trái tim tôi siết lại. Anh chị\dots\ Onii-sama và Onee-sama. Tôi nuốt khan, bật ra câu hỏi run rẩy: ``\eng{Onii-sama\dots\ và Onee-sama\dots\ họ có\dots\ có ổn không?}''

Bên kia ngập ngừng. Rồi giọng nói ấy đáp, ngắn gọn: ``\eng{Anh trai của cô thì ổn. Nhưng chị của cô thì\dots}''

Âm thanh lạc đi, như bị nuốt trọn bởi chính bầu không khí lạnh lẽo nơi đây. Tôi lảo đảo, ngực thắt lại, đầu óc choáng váng. Trái tim tôi gõ loạn như sắp vỡ tung.

``\eng{T-Tôi tin là chị của cô sẽ ổn thôi, tôi sẽ cố gắng để liên lạc được với cô ấy. Bảo trọng!}'' giọng nói khẽ dịu lại, rồi biến mất hoàn toàn, để lại tôi một mình trong hành lang chập chờn ánh sáng và tiếng lạch cạch của cửa gỗ đóng mở.

Tôi đứng chết lặng. Bàn tay run rẩy cầm đèn pin như muốn rơi xuống. Nỗi sợ hãi, tuyệt vọng và trống rỗng tràn đến, nuốt chửng tôi. Nếu Onee-sama gặp chuyện gì\dots\ nếu chị ấy thật sự biến mất\dots\ tôi phải sống thế nào đây? Một dòng nước mắt nóng hổi lăn dài xuống má, rơi vào khoảng tối lặng im.

Nhưng\dots\ tôi lại nhớ đến ánh mắt chị, đến bàn tay dịu dàng nhưng mạnh mẽ của chị, từng dạy tôi rằng: ``Reiko, nếu em chỉ biết khóc và sợ hãi, em sẽ mãi mãi chẳng thể nào trưởng thành được đâu.''

Tôi hít một hơi thật sâu, gạt mạnh nước mắt, rồi ngẩng đầu lên. Dù trong lòng tôi vẫn run sợ, nhưng tôi biết mình không thể đứng đây mãi. Tôi phải đi tiếp, phải chứng minh rằng tôi không vô dụng. Rằng tôi có thể tự bước, tự tìm con đường của riêng mình, ngay cả khi bóng tối đang chực chờ nuốt trọn.

Tôi giơ cao đèn pin, ánh sáng hắt vào hành lang dài hun hút phía trước. Những cánh cửa vẫn đóng mở liên hồi, như đang thách thức. Tôi cắn chặt môi.

``Nếu đây là thử thách\dots'' tôi thì thầm với chính mình, ``\dots thì tôi sẽ đối mặt. Dù sợ hãi, tôi cũng sẽ bước qua.''

Và thế là, với đôi chân run rẩy nhưng kiên định, tôi lại tiến lên.

Bằng tất cả sự yếu đuối và mạnh mẽ đan xen trong trái tim này, tôi bước tiếp về phía hành lang hun hút, nơi mỗi nhịp thở đều nặng nề như sắt nung đỏ.

Tôi hít sâu một hơi, bước từng bước dọc hành lang. Những tiếng lạch cạch! của cửa trượt cứ dội vào tai, dồn dập như muốn xé toang màng nhĩ. Tim tôi đập loạn nhịp, nhưng tôi vẫn bước. Tôi biết, nếu dừng lại, nếu quay đầu, tôi sẽ chẳng bao giờ thoát ra nổi.

Ánh sáng từ đèn pin run rẩy quét qua một lớp học. Và ngay khoảnh khắc ấy, tôi suýt đánh rơi nó.

``\textbf{Eeek!}''

Bên trong là những bóng người. Không, là những cái bóng đen mờ ảo, méo mó, dáng dấp quen thuộc đến mức tôi thấy nghẹn thở. Tôi nhận ra Nanase-san, với kiểu đứng quý phái đặc trưng. Yae-san, với thân hình vạm vỡ. Shiina-san, tỏ vẻ du côn lạnh lùng như thường lệ. Và cả Honoka-san, đôi vai nhỏ bé như đang run rẩy trong sương mù.

``Không\dots\ thể nào\dots'' tôi lùi lại một bước, nhưng cánh cửa sau lưng rầm một tiếng đóng sập, làm tôi giật mình ngã xuống. Trong lớp, bàn ghế bắt đầu bay lên, xoay vòng lơ lửng, như bị một sức mạnh vô hình giật dây. Tiếng gỗ cọt kẹt, tiếng kim loại loảng xoảng, tất cả hợp lại thành một thứ âm thanh ma quái khiến tôi muốn bịt tai mà gào lên.

Trong đầu tôi chợt hiện lại giọng Shino-san, vang vọng, dịu dàng nhưng nghiêm khắc: ``Đừng chạm vào họ. Đừng để bản thân bị cuốn theo. Những bóng đó không còn là con người nữa.''

Rồi là giọng của Onii-sama, trầm chắc như trụ cột: ``Nếu em sợ, hãy nhớ rằng đôi chân em vẫn còn. Chúng sẽ đưa em đi qua nỗi sợ.''

Và cả Onee-sama, vừa nghiêm vừa dịu: ``Em có thể khóc, có thể run, nhưng tuyệt đối không được dừng lại. Nếu ngã xuống, hãy đứng lên.''

Tôi cắn môi, cảm giác mặn chát của máu lan ra, kéo tôi khỏi nỗi choáng váng. ``\textit{Đúng\dots\ mình không được tiếp xúc với họ. Không được để chúng kéo mình vào.}'' tôi lẩm bẩm như một câu thần chú.

Bàn ghế bắt đầu lao ra khỏi lớp, đập mạnh xuống sàn, khiến hành lang rung lên từng hồi. Tôi nép người vào tường, mồ hôi lạnh túa ra trên trán, nhưng tôi vẫn giữ chặt đèn pin, đưa ánh sáng soi về phía trước. Những bóng đen trong lớp bắt đầu xoay người, dường như đang dõi theo tôi, đôi mắt vô hình bám riết lấy từng bước chân tôi.

``Không nhìn\dots\ không nghe\dots\ không để ý\dots'' Tôi tự thì thầm, bước thật nhanh qua những cánh cửa đang rung lắc, mặc cho mỗi tiếng đóng mở như búa giáng xuống màng nhĩ.

Một cánh cửa bật mở mạnh hơn hẳn, gió lạnh phả vào mặt tôi. Trong tích tắc, tôi thấy bóng Nanase-san vươn tay ra, những ngón tay dài ngoằng, run rẩy, gần như chạm đến má tôi. Tôi thót tim, nhưng ngay khoảnh khắc ấy, tôi tránh kịp sang một bên.

Tôi nhắm chặt mắt, quay đi, ép bản thân bước tiếp. Đèn pin hắt sáng lên sàn gỗ loang lổ, dẫn đường cho tôi. Từng bước, từng bước một, dù run rẩy nhưng không dừng lại.

Tiếng la hét, tiếng bàn ghế vỡ tung, tất cả như đang gào vào tai tôi rằng: ``\textbf{Đứng lại! Quay lại! Đừng đi nữa!}'' Nhưng tôi biết, nếu tôi nghe theo, tôi sẽ bị chôn vùi mãi mãi trong quá khứ này.

Tôi nghiến răng, vừa chạy vừa tự nhủ với bản thân: ``\textit{Mình có thể làm được\dots\ Mình có thể làm được! Mình phải rời được khỏi đây!}''

Vừa chạy được vài bước, toàn bộ hành lang bỗng nổi sóng. Một trận động đất ập đến chỉ trong vài giây, nhưng đủ khiến tường gạch rên rỉ, xà ngang kêu cót két, sàn gỗ gập ghềnh nhấc lên rồi hạ xuống. Tôi loạng choạng, đầu gối khuỵu, đèn pin văng khỏi tay, lăn lông lốc dưới sàn. Ánh sáng loang loáng quét qua khe cửa, vẽ những vệt sáng xé toang bóng tối rồi tắt ngóm. Tôi nghe rõ tiếng gạch vỡ, tiếng sắt rít, tiếng búa đập vào xương sống mình. Một mảnh thạch cao rơi xuống cách đầu tôi gang tấc, bụi phủ mờ mắt. Tôi ho sặc sụa, mùi bụi, mùi máu, mùi mồ hôi hòa lẫn.

Nhưng tôi không dám nằm im.

Tôi lê người về phía đèn pin, tay run rẩy vớ lấy, dập lại vào tường để bật lại ánh sáng. Chùm sáng yếu ớt le lói, run rẩy như trái tim tôi. Tôi gào thét trong đầu: ``\textit{Đứng dậy! Đứng dậy ngay!}'' Đôi chân tôi bật lên, loạng choạng bám vào cửa sổ vỡ, bám vào gờ tường nứt, bám vào bất cứ thứ gì không sụp đổ. Tôi lê từng bước như người say, nhưng không ngừng.

Trận động đất qua đi nhanh như lúc đến, chỉ để lại hành lang rúng động, bụi mù mịt, và trong lòng tôi một ngọn lửa nhỏ vừa bị dập tắt lại bùng cháy dữ dội hơn. Tôi ôm chặt đèn pin, tiếp tục chạy, máu từ khóe môi rỉ xuống cằm, nhưng tôi không lau. Tôi chỉ cần tiến về phía trước.

Và cuối cùng, như xuyên qua một bức màn vô hình, tiếng động dần lắng xuống. Những cánh cửa ngừng rung. Khi tôi mở mắt, hành lang phía trước trống trải, yên lặng đến mức tôi nghe rõ tiếng thở hổn hển của chính mình.

Tôi khuỵu xuống, ôm ngực, cảm giác như vừa thoát chết trong gang tấc. Nhưng một ngọn lửa nhỏ đã nhen lên trong lòng tôi. Tôi đã vượt qua thử thách đó.

\ct

Tôi bước chậm trên sàn gỗ cũ, ánh sáng đèn pin chiếu ra phía trước, nhịp tim vẫn chưa kịp ổn định. Giữa hành lang, một chiếc ghế gỗ cũ kỹ nằm chỏng chơ, nghiêng lệch như thể vừa có ai đó ngồi lên một lát rồi biến mất. Tôi ngần ngại tiến lại gần, ánh sáng đèn pin phản chiếu lên lớp bụi mỏng phủ trên mặt ghế.

Rồi\dots\ một bóng người xuất hiện.

Tôi giật mình, suýt đánh rơi đèn. Ngay trước mắt, một cô gái tóc nâu-đỏ, cột hai bím thấp, dáng nhỏ nhắn - thật lạ là mái tóc ấy, gương mặt ấy, đâu đó\dots\ lại giống tôi. Cô ngồi lặng trước khung tranh dựng nghiêng, trong tay là một cây cọ, nhưng bàn tay run rẩy, không hề di chuyển.

``Bức tranh này không được rồi\dots''
giọng cô ta run rẩy, khản đặc. ``Không được! Mình chưa bao giờ cảm thấy bí như vậy cả. Tại sao mình lại không vẽ được chứ\dots''

Tôi đứng chết lặng. Trong đầu tôi, một mẩu chuyện vụn vỡ hiện lên: ``Lời nguyền của phòng mỹ thuật.'' Một cô gái muốn tạo ra kiệt tác để lưu danh, nhưng khi không thể, cô đã dùng chính máu của mình để hoàn thành bức tranh cuối cùng. Sau đó, cô đã phải trả giá bằng chính mạng sống của mình.

``\dots Chẳng lẽ, là cô ấy sao?''

Cô gái nâu-đỏ đứng dậy, ngửa mặt nhìn lên trần gỗ, đôi mắt hai màu vàng-đỏ trống rỗng. ``Tại sao lại như vậy chứ\dots\ Mình biết mình muốn vẽ mà\dots\ nhưng mà bàn tay của mình không chịu nghe mình cả\dots\ Mình không hiểu nổi. Mình không thể chịu nổi nữa\dots''

Giọng nói ấy vang vọng, nghẹn ngào, rồi đột ngột lịm đi. Ánh sáng đèn pin của tôi chập chờn, và khi tôi mở mắt lại, chỉ còn chiếc ghế chỏng chơ như ban đầu. Không có khung tranh, không có người.

Ngoại trừ việc, một bức ảnh nhỏ nằm trên mặt ghế, nơi mà cô ta vừa ở đó.

Tôi nhặt nó lên.

Trong bức ảnh là ba người: một chàng trai tóc tổ quạ, mặc áo phông đen; một cô gái tóc đuôi ngựa xù, áo trắng, váy đỏ; và một cô gái tóc hai bím thấp, áo vàng, váy xanh, là tôi.

Dưới khung ảnh là ba cái tên, được viết rõ ràng:

``Hikawa Kuon.''

``Hikawa Kaigou.''

``Hikawa Suzuki.''

Tôi nhìn ba khuôn mặt ấy. Cảm giác quen thuộc ập đến. Onii-sama và Onee-sama. Nhưng tại sao\dots\ tại sao tên lại khác như vậy? Giống hệt như cái cuốn kỷ yếu cũ mà tôi từng thấy ở trên tàu trong vòng lặp trước.

Tôi đưa tay run rẩy chạm vào cái tên cuối cùng, Hikawa Suzuki.

``!?''

Ngay khoảnh khắc ấy, một cơn đau dữ dội như xé đầu tôi ra. Tôi hét khẽ, gục xuống, hai tay ôm đầu. Mọi thứ như một cơn lũ dữ, hình ảnh, âm thanh, ký ức, tràn ngập, va đập, quăng quật tôi giữa dòng thời gian.

Tôi thấy một ngôi nhà nhỏ ngập ánh nắng, tiếng cười vang lên.

Thấy bố mẹ tôi đang nở nụ cười hạnh phúc\dots\ rồi tiếng phanh xe rít lên, tiếng kính vỡ\dots\ và bóng tối.

Tôi thấy Onee-chan - Kaigou - ôm tôi khóc, run rẩy, thề rằng sẽ nuôi tôi lớn dù chỉ còn hai chị em.

Tôi thấy bảy năm trôi qua trong cô độc nhưng đầy ấm áp.

Rồi\dots\ một đêm nọ. Tiếng nổ. Ngọn lửa. Ngôi nhà bốc cháy.

Tôi và chị chạy, nhưng ngọn lửa nuốt chửng tất cả.

Ánh sáng trắng, và The Void.

Tôi thấy chính mình và chị bị tách ra giữa khoảng không đen vô tận.

Tôi đã nghĩ mình sẽ biến mất mãi mãi\dots cho tới khi anh ấy xuất hiện.

Một bàn tay vươn tới, kéo tôi ra khỏi khoảng không đó.

Hikawa Kuon.

Onii-sama\dots\ người đã đưa tôi đoàn tụ với Onee-sama.

Từ đó, ba chúng tôi cùng nhau bước qua những thế giới lạ lẫm.

Tôi thấy ``There Is No Game,'' thấy Game-san lần đầu tiên hiện lên trên chiếc đồng hồ, và chính tôi là người đã thuyết phục ông ta đi cùng.

Rồi bốn người chúng tôi - gồm tôi, chị Kaigou, anh Kuon, và Game-san - lạc vào thế giới Nijisanji. Tôi thấy lại gương mặt của Chihiro-chan và Ririmu-han, thấy họ cười, giơ máy ảnh, nói: ``Cười lên nào, ba người! Tấm ảnh này sẽ làm kỷ niệm đấy!''

Và\dots\ đó chính là bức ảnh này.

Sau đó, bầu trời xé toang. Một lỗ xoáy khổng lồ mở ra, cuốn tất cả vào trong.

Tôi thấy ánh sáng trắng, rồi không còn gì nữa.

Tất cả vỡ vụn, rồi chìm vào tĩnh lặng.

Khi tôi mở mắt, hơi thở nặng nề, mồ hôi lạnh ướt trán. Bức ảnh vẫn nằm trong tay tôi, nhưng trái tim tôi đã không còn trống rỗng như trước.

Tôi siết chặt nó, nước mắt lăn dài trên má.

``\dots Mình nhớ được rồi.''

Giọng tôi khản đặc, nhưng dứt khoát. ``Tôi không phải là Haragen Reiko. Tôi là \emph{Hikawa Suzuki}. Là tôi, em gái của Onee-sama - Hikawa Kaigou, và là nhận Onii-sama - Hikawa Kuon - là anh.''

Một làn gió nhẹ thổi qua hành lang, như một lời đáp mơ hồ.

Tôi lau nước mắt, hít sâu, rồi đứng thẳng dậy.

``Onii-sama, Onee-sama\dots\ em sẽ tìm được hai người. Dù có phải đi qua bao nhiêu thế giới, em cũng sẽ không bỏ cuộc.''

Tôi khẽ khép mắt lại. Cơn đau nhói trong đầu vẫn còn đó, nhưng lần này\dots\ nó mang lại cảm giác khác. Không còn là nỗi sợ hay hỗn loạn như trước, mà là một thứ cảm xúc tràn đầy: sự ấm áp, thân quen, và cả\dots\ nhớ nhung.

Hơi thở tôi nghẹn lại khi nghĩ về họ. Hình ảnh hai người - Onii-sama với nụ cười nhẹ mà luôn giấu đi mệt mỏi, Onee-sama với ánh mắt kiên định và cái vuốt tóc đầy ân cần - hiện lên rõ ràng trong tâm trí tôi. Tôi muốn gặp họ. Tôi phải gặp lại họ. Dù là ở đâu, dù phải vượt qua bao nhiêu thứ đáng sợ nữa, tôi cũng không thể dừng lại.

Tôi liếc nhìn chỗ chiếc ghế khi nãy. Cô gái tóc nâu-đỏ đã biến mất rồi, chỉ còn lại bức ảnh tôi vẫn nắm chặt trong tay. Vẻ tuyệt vọng của cô ấy vẫn còn vương lại trong tim tôi, như một lời nhắc: nếu mình đánh mất niềm tin, mình cũng sẽ tan biến như cô ấy thôi.

Tôi cúi đầu, nói nhỏ như tự nhủ với bản thân: ``Mình sẽ không như vậy đâu\dots\ Cảm ơn nhé.''

Tôi đứng dậy, tay vẫn nắm bức ảnh, hướng tới cánh cửa trượt gỗ phía cuối hành lang. Bên kia, một thứ ánh sáng trắng rực rỡ xuyên qua khe cửa, sáng đến mức khiến tôi phải nheo mắt. Ánh sáng ấy vừa đẹp, vừa đáng sợ, như thể tôi chỉ cần bước qua thôi là sẽ bị nuốt chửng hoàn toàn.

Tôi do dự một chút. Nhưng rồi, tiếng nói trầm khàn khi nãy của Game-san - ``Hãy tự tin vào bản thân'' - lại vang vọng trong tâm trí.
Tôi hít một hơi sâu, nắm chặt tay, và bước về phía trước.

Khi bàn chân tôi vượt qua mép cửa, tất cả mọi thứ xung quanh lập tức hóa thành trắng xóa. Tôi không còn cảm thấy mặt đất dưới chân nữa, chỉ có thứ ánh sáng nhấn chìm toàn thân. Nhưng lần này, tôi không sợ. Tôi biết\dots\ nơi bên kia, Onii-sama và Onee-sama đang chờ.

\dots

\dots

\pov{???}

\dots

``Vậy\dots\ hai người có vẻ như là các học sinh mới chuyển đến à.''

Giọng nữ ấy vang lên giữa khoảng không mơ hồ, nhẹ mà rõ. Tôi biết giọng đó. Shino. Tôi đã từng nghe nó, một âm điệu đều đều, lạnh mà trầm, như thể người nói luôn nhìn thấu mọi thứ.

Rồi một tiếng cười khe khẽ vang lên. Nó nghe có gì đó giễu cợt, xa xăm khỏi tôi.

\dots Tôi vẫn còn sống\dots ư?

\dots

\dots

``Gah!''

Tôi bật mở mắt. Một luồng ánh sáng mờ chạy ngang qua, lướt nhanh, rồi biến mất. Khi đôi mắt tôi dần thích nghi, cảnh tượng hiện ra khiến tôi ngạc nhiên.

Tôi đang ở trong một toa tàu. Những hàng ghế trống im lặng, cửa kính phản chiếu gương mặt tôi nhợt nhạt. Trông rất giống với con tàu mà chúng tôi đã phanh. Bên ngoài, bóng tối trôi ngược lại với vận tốc đáng sợ, như thể con tàu đang lao vào hư vô.

Tôi đang ở trên một con tàu khác, không phải con tàu mà ba chúng tôi đã phanh gấp hồi nãy. Tiếng bánh xe chạy trên đường ray hiện hữu, băng băng đi về phía trước như chưa có gì xảy ra.

``\dots Nii-chan? Reiko?'' tôi gọi khẽ, nhưng chỉ có tiếng bánh sắt nghiến trên đường ray đáp lại. Không có một giọng nào khác.

Tôi phải mất vài giây mới có thể ngồi dậy, tay bấu vào tay vịn lạnh toát. Không có Kuon bên cạnh. Không có Reiko. Chỉ có tôi, một mình, trong không gian im lặng chết chóc này.

Tôi quay nhìn khắp nơi, hy vọng có một bóng người nào đó, nhưng chẳng có gì ngoài hàng ghế trống dài vô tận và đèn trong đường hầm vẫn chạy đi chạy lại như không có điểm dừng. Ngay cả âm thanh của Shino-san cũng biến mất rồi, như thể nó chưa từng tồn tại.

Một cảm giác lạnh buốt len qua sống lưng. Tôi ôm lấy tay mình, cố giữ bình tĩnh, nhưng càng lúc tim tôi càng đập nhanh hơn. Không có ai ở đây. Không có gì quen thuộc.

Chỉ có tôi và con tàu này, đang đưa tôi đến nơi mà tôi chưa từng biết.

``Tôi\dots đang ở đâu thế này?''

*XẸT XẸT*

Tiếng rè rè chát chúa từ loa phát thanh trên đầu. Tôi giật mình, ngẩng lên.

``\eng{Thông báo, thông báo\dots\ Hikawa Kaigou-san có ở đây không?}''

Giọng nói ấy khàn, hơi đục, và mang một âm điệu lạ lùng, đậm chất Pháp, lại còn nói bằng tiếng Anh. Nó không giống giọng nói của một người điều khiển tàu, mà giống như một thứ gì đó khá giống máy móc, nhưng lại rất\dots\ con người.

Tôi chau mày, cố lắng nghe cho rõ. ``Hikawa\dots\ Kaigou?''

Cái tên ấy khiến tôi thoáng khựng lại. Tôi chưa từng nghe thấy nó bao giờ, hoặc ít nhất là tôi nghĩ vậy. Nhưng\dots\ tại sao tôi cảm thấy tim mình dồn dập như vậy?

Rồi, như một dòng điện chạy qua đầu, tôi nhớ đến cuốn kỷ yếu mà tôi từng thấy trên con tàu trong vòng lặp trước. Cái tên đó\dots\ đã từng nằm ở đó. Cái tên trên tấm ảnh có khuôn mặt rất giống tôi.

``\eng{\dots T-Tôi á?}'' tôi chỉ tay vào ngực mình, gần như bật ra câu hỏi theo phản xạ.

``\eng{Err\dots\ Đúng. Cô đấy.}'' giọng nói ấy đáp lại, khẽ pha chút tiếng cười mệt mỏi. ``\eng{\hl{Có vẻ như cô cũng bị chưa nhớ ra được bản thân rồi.}}''

Tôi im lặng, không biết phải nói gì. Giọng nói đó tiếp tục, trầm nhưng rõ: ``\eng{Tôi không có nhiều thời gian đâu, nên nói ngắn gọn thôi: Đừng tin vào những gì bản thân mình thấy. Ngoại trừ cái ánh sáng trắng.}''

``\eng{Ánh sáng trắng?}'' tôi nhắc lại, cố hiểu xem ông ta muốn nói gì, nhưng ông đã nói trước:

``\eng{Ừ. Còn về việc cô đang ở đâu thì\dots\ phức tạp lắm. Có nói cô cũng không hiểu đâu. Cứ hiểu nôm na là cô đang ở trong một nơi hoàn toàn khác.}''

Tôi cau mày, cố gắng nắm bắt từng lời. ``Một nơi khác\dots'' câu đó khiến tôi lạnh sống lưng.

Thế rồi, tôi chợt nhớ đến hai anh em của tôi.``\eng{Vậy\dots\ còn hai người kia thì sao? Nii-chan, và Reiko?}''

Có một thoáng im lặng. Rồi giọng nói kia đáp, có vẻ như đang kiểm tra dữ liệu gì đó: ``\eng{Nii-chan? À, ý cô là cậu con trai có ngoại hình giống cô ư? Có, cậu ấy vẫn ổn.}''

Tôi khẽ thở phào, nhưng rồi ông ta nói tiếp: ``\eng{Còn Reiko của cô thì\dots\ cô ấy đang ở một không gian khác. Không cùng nơi với cô đâu.}''

Tôi nín thở. Cổ họng khô rát lại. Một không gian khác\dots\ có nghĩa là chúng tôi lại tách biệt. Câu đó thôi đã khiến tôi rùng mình. Tôi đã cảm tưởng rằng đã từng mất Reiko một lần rồi, trong một khoảnh khắc mờ ảo mà tôi chẳng thể nắm bắt. Tôi\dots\ không muốn điều đó lặp lại nữa.

Dường như ông ta biết tôi đang nghĩ gì, vì ngay sau đó, giọng nói ấy dịu đi, có phần trấn an: ``\eng{Tôi tin là cô ấy sẽ ổn thôi.}''

Tôi siết chặt nắm tay, nhìn xuống sàn tàu. Cổ họng tôi như thể bị thắt lại, vừa vì sợ, vừa vì một nỗi hối tiếc mơ hồ. Tôi đã không thể bảo vệ nó lần trước. Tôi sẽ không để điều đó tái diễn.

``\eng{Khoan đã--}'' tôi toan hỏi thêm, nhưng giọng nói kia đã ngắt lời: ``\eng{Không còn thời gian đâu. Chúc cô may mắn.}''

Ngay sau đó, tiếng rè rè vang lên một lần nữa, kéo dài như một hơi thở cuối, rồi im bặt. Chỉ còn tiếng lăn bánh nặng nề của đoàn tàu giữa khoảng không vô định.

Tôi ngồi lặng, nghe tim mình đập mạnh trong lồng ngực. Những cái tên ``Kaigou'', ``Nii-chan'', ``Reiko'', tất cả bắt đầu xoáy trong tâm trí tôi như những mảnh ký ức đang tìm cách ghép lại. Một nỗi sợ mơ hồ trỗi dậy, rằng sẽ lại đánh mất cô em gái mà tôi đã hứa sẽ bảo vệ đến cùng.

Tôi hít sâu, đứng dậy. Mắt nhìn ra ô cửa kính đen đặc phía ngoài.

``Trước mắt thì\dots\ phải tìm cách thoát khỏi đây đã.''

Dù con tàu này có đưa tôi đến đâu, tôi phải tìm được họ.

``Cả hai người\dots\ làm ơn hãy bình an vô sự\dots''

Tôi cúi xuống, nhặt lại chiếc đèn pin nhỏ mà ai đó - có lẽ là chính tôi trong vòng lặp trước - đã đánh rơi. Ánh sáng vàng nhạt chiếu rọi về phía trước, quét dọc theo sàn tàu phủ đầy bụi. Tiếng kim loại va vào nhau theo nhịp rung của đường ray khiến tôi bất giác siết chặt cán đèn pin hơn. Toa tàu dài, và ánh sáng ấy, dù yếu ớt, vẫn đủ để khiến không gian trước mắt tôi trở nên\dots\ chân thực đến đáng sợ.

Tôi bước chậm lại khi thấy ở chỗ ngoofi cuối toa có hai bóng dáng. Cả hai ngồi đối diện nhau, ngay bên cửa nối toa, mang dáng dấp Akane-san và Yuka-san. Tôi nhận ra họ ngay, từ mái tóc, bờm tai thú, cho tới bộ đồng phục mà họ mặc lần cuối.

Nhưng, họ không nói gì cả. Chỉ lặng lẽ nhìn nhau, y như những bức tượng người đã bị thời gian nuốt trôi.

Tôi khựng lại. Cổ họng nghẹn cứng, và trong đầu chợt vang lên những hình ảnh cũ.

Ngày ấy, khi lần đầu gặp họ, tôi đã nghĩ họ chỉ là một cặp bạn thân đơn thuần, lúc nào cũng cười nói, cũng chọc nhau bằng những chuyện vặt.

Và rồi, khi Yuka biến mất. Tôi còn nhớ ánh mắt của Akane lúc đó. Tôi dám cá là cô ấy hoảng loạn, gọi tên bạn mình đến khản cả giọng, điên cuồng chạy qua từng dãy hành lang, chẳng màng đến những lời can ngăn.

Cảnh tượng ấy\dots\ tôi hiểu hơn ai hết. Vì tôi cũng đã từng như thế, khi Reiko biến mất giữa một vòng xoáy vô định, tôi cũng đã mất lý trí đến mức sẵn sàng đánh đổi mọi thứ chỉ để kéo em ấy trở lại.

Nhưng rồi, khi tôi chớp mắt, cả Akane lẫn Yuka đều tan biến, như thể chưa từng tồn tại. Ánh sáng từ đèn pin hắt lên, chiếu xuyên qua khoảng trống mà họ để lại. Một cảm giác trống rỗng lan ra trong lồng ngực. Cảm giác khi thứ mà mình tưởng đã chạm tới lại hóa ra chỉ là ảo ảnh.

Tôi hít một hơi thật sâu.

``\dots Đúng, chỉ là ảo giác thôi,'' tôi tự nhủ với bản thân.

Tôi tiếp tục bước tới. Cánh cửa nối toa mở sẵn, lặng lẽ đón tôi như thể đang mời gọi.

\ct

Toa kế tiếp yên tĩnh đến lạ. Không còn những bộ bàn ghế chắn ngang đường đi, giống như ở vòng lặp trước. Đèn trần nhấp nháy từng hồi, ánh sáng chập chờn như đang cố thở. Không khí loãng và lạnh, đủ khiến tôi rùng mình.

Tôi bước từng bước, lắng nghe tiếng giày mình vang lên trên sàn tàu rỗng.

Nhưng khi tôi ngẩng đầu, ánh đèn pin bất chợt chiếu tới một khung cảnh khiến tôi chết lặng. Ở phía xa là Yuka-san. Vẫn dáng đi đó, vẫn mái tóc ấy\dots\ nhưng cô không đi một mình. Cô đang được ai đó dẫn đi. Người ấy có dáng dấp của Shino. Mái tóc nâu dài kia, cách bước đi chậm rãi đó\dots\ tôi không thể nào nhầm được.

Cổ họng tôi khô khốc. Cảnh tượng ấy\dots\ tôi đã từng thấy rồi. Trong giấc mơ, giữa vòng lặp cũ, tôi từng thấy Yuka đi theo ai đó giống hệt Shino, như thể bị thôi miên, vô thức, trống rỗng. Và giờ, cảnh đó lại đang tái diễn trước mắt tôi, nhưng rõ ràng hơn, thật hơn, đến mức tôi có thể nghe thấy tiếng giày của họ chạm xuống sàn.

``\textbf{Yuka-san!}'' tôi gọi lớn. Nhưng, tôi cảm tưởng bên kia không nghe được những gì. Chỉ có tiếng đèn nhấp nháy và tiếng gió lạnh trườn qua.

Khi tôi cố bước nhanh hơn, cả hai bóng dáng ấy đột nhiên biến mất ngay trước cửa nối toa. Như thể mất tích không một dấu vết.

Một phần trong tôi muốn gào lên, muốn chạy theo, muốn kéo cô ấy ra khỏi cái gì đó đang thao túng cô. Nhưng phần còn lại hiểu rằng,nếu tôi bước thêm một bước, tôi có thể sẽ không bao giờ quay lại được.

Chưa kể rằng, tôi cũng đã biết số phận của cô ấy cùng cô bạn thân và Uzuki-san rồi: họ đã hoàn toàn biến mất khỏi thế gian này.

Tôi cúi đầu, khẽ cười. ``Vẫn là trò thử lòng, đúng không?''

Rồi tôi siết chặt đèn pin, bước thẳng tới. Cánh cửa tới toa tiếp theo lại được mở sẵn.

\ct

Toa tàu tiếp theo tối om khi tôi bước vào. Vừa qua khỏi ngưỡng cửa, ánh sáng từ chiếc đèn pin yếu ớt của tôi phản chiếu lên những hàng ghế dài, những ô cửa sổ phủ đầy bụi, và rồi toàn bộ đèn trần phụt tắt.

Một luồng gió lạnh lùa qua, khiến tôi rùng mình. Rồi\dots\ *BỊCH BỊCH BỊCH*, từ bên ngoài, hàng ngàn cánh tay bắt đầu đập mạnh vào khung kính cửa sổ, dồn dập, liên hồi, như thể có cả một biển người đang cố xông vào toa tàu. Tiếng va đập kim loại, tiếng kính rung, tiếng móng tay cào trượt để lại vệt dài trắng xoá\dots\ tất cả hòa thành một thứ âm thanh hỗn loạn đến chói tai.

Tôi rọi đèn pin về phía một cửa sổ. Ánh sáng lướt qua những hình thù mờ ảo, những cánh tay gầy guộc, biến dạng, quờ quạng đòi được chạm vào tôi. Trông như thể thế giới bên ngoài đang sống dậy, nhưng lại là một thế giới không còn hình người.

Tôi bước chậm dọc hành lang, cố giữ bình tĩnh. Tiếng thở của mình vang lên nặng nề giữa âm thanh đập cửa dồn dập kia. Bất chợt\dots

*KÍÍÍÍÍÍÍÍÍÍÍT!!!*

\dots con tàu phanh gấp, cả cơ thể tôi bị hất về phía trước. Mọi thứ nghiêng ngả. Tôi lảo đảo, khuỵu gối xuống sàn, chống tay lại theo phản xạ. Cổ tay đau rát, miệng bật ra một tiếng rên khe khẽ.

``\dots Agh\dots\ cái quái gì vậy\dots''

Tôi nằm đó vài giây, thở gấp. Cơn choáng chóng mặt vẫn chưa qua, nhưng khi tôi ngẩng đầu dậy, ánh đèn pin lướt qua một vật lạ. Ở hàng ghế thứ ba tính từ cuối toa, có một thứ lấp lánh phản chiếu ánh sáng. Đó là một cái nơ buộc tóc cầu vồng, với những dải ruy-băng sặc sỡ như vừa được cắt ra từ một chiếc hộp quà Giáng Sinh. Màu đỏ, xanh, vàng, tím đan xen nhau, và cách buộc thì\dots\ thật sự lòe loẹt.

Tôi cau mày. ``Trông thứ này\dots\ mình đã thấy ở đâu rồi\dots''

Ngay khi những từ đó rời khỏi môi, đầu tôi như thể bị ai đó giáng mạnh bằng búa. Một cơn đau buốt xé toạc hai thái dương. Tôi ôm đầu, gục xuống sàn, hơi thở đứt quãng.

Và rồi, ký ức tràn về như một con sóng dữ.

Tôi thấy mình - Kaigou, đứa con gái lớn trong gia đình Hikawa - cùng một cô em gái nhỏ có đôi mắt sáng trong, luôn gọi tôi bằng giọng dịu dàng: ``Onee-sama ơi\dots''

Hikawa Suzuki.

Đúng vậy\dots\ đó là em gái tôi.

Cả hai chúng tôi đã mất cha mẹ trong một tai nạn, trở thành trẻ mồ côi suốt bảy năm. Tôi đã làm tất cả để bảo vệ em, vừa là chị, vừa là mẹ, vừa là người che chở duy nhất còn lại. Mọi người gọi tôi là ``cực đoan'' vì cách tôi bảo vệ em khỏi những kẻ đàn ông bệnh hoạn ngoài xã hội thối nát kia, nhưng tôi chẳng quan tâm. Nếu thế giới này chỉ còn lại một người, thì tôi muốn người đó là em.

Và rồi, đêm định mệnh ấy ập tới.

Lửa.

Tiếng nổ.

Tiếng em khóc gọi tôi giữa khói đặc quánh.

Chúng tôi chạy, nắm chặt tay nhau, nhưng rồi ánh sáng đỏ ấy nuốt chửng tất cả, và khi tôi mở mắt ra, tôi thấy mình lạc ở một nơi\dots\ The Void.

Một khoảng không chỉ toàn màu đen, vô tận, lạnh đến tê tái.

Tôi và em đã bị tách khỏi nhau. Tôi gào khản cổ, chạy trong vô định trong suốt thời gian mà tôi không tài nào nhớ nổi. Tôi đã tưởng rằng tôi đã lạc mất con bé mãi mãi\dots

\dots\ cho đến khi một người con trai xuất hiện, đôi mắt đỏ sắc lẹm, giọng nói bình tĩnh và dứt khoát.

Hikawa Kuon. Tôi gọi cậu ta là ``Kuon-kun.''

Cậu ấy đã giúp tôi và em đoàn tụ, đã hứa sẽ đưa chúng tôi đến một cuộc sống mới. Tôi còn nhớ mình đã lạnh lùng, nghi ngờ, vì tôi không tin đàn ông. Nhưng rồi, chính cậu ta đã chăm sóc Suzuki trong khoảng thời gian tìm tôi\dots\ và tôi hiểu, cậu ta không nói dối.

Sau khi thoát khỏi The Void, chúng tôi đi qua nhiều thế giới. ``There Is No Game'', rồi thế giới của Nijisanji. Từng ngày, từng khoảnh khắc, tôi cảm giác như mọi thứ dần trở lại bình thường. Tôi còn nhớ ngày thứ ba ở đó, khi Suzuki và Kuon-kun cùng làm món đồ thủ công kỳ lạ: một chiếc nơ cầu vồng. Cậu ta từng nhận xét nó là ``hơi lòe loẹt và ố dề,'' tôi chỉ bật cười nhạt, nhưng trong lòng, tôi biết đó là món quà đẹp nhất mà tôi từng nhận.

Rồi, lỗ xoáy.

Cơn gió cuộn lên, hút tất cả bọn tôi vào. Tôi thấy ánh sáng trắng nuốt trọn mọi thứ.

Từ đó, tôi không còn nhớ gì nữa.

Và giờ đây, trong toa tàu đen đặc này, khi tôi cầm chiếc nơ cầu vồng trên tay, những ký ức ấy trở về nguyên vẹn.

Tôi đứng lặng. Trái tim đập loạn trong lồng ngực. Mắt tôi nhòe đi, không rõ vì nước mắt hay vì ánh đèn pin quá chói.

``\dots Mình là\dots\ \emph{Hikawa Kaigou}\dots'' tôi thốt lên trong hơi thở run rẩy.

Tôi nhớ ra rồi. Tôi không phải Haragen Saikyou, không phải một bóng người lạc trong mớ ký ức nhân tạo. Tôi là chính mình, người chị đã thề sẽ bảo vệ em gái mình bằng cả mạng sống.

Tôi khẽ mỉm cười. Một nụ cười vừa nhẹ nhõm, vừa đau lòng.

``\dots Cuối cùng thì\dots\ chị cũng nhớ ra rồi, Suzu\dots''

Ngoài kia, tiếng đập cửa sổ vẫn vang lên, nhưng lần này, tôi không còn sợ nữa.

Tôi siết chặt chiếc nơ cầu vồng trong tay, thứ duy nhất nối tôi với những người thân yêu, rồi ngẩng đầu lên, nhìn về phía cánh cửa nối toa vẫn đang hé mở.

Phía sau đó\dots\ là ánh sáng mờ nhạt, yếu ớt, nhưng thật đến lạ thường.

\dots

\dots

\pov{Hikawa Kuon}

Tôi đứng im lặng giữa bên trong nhà ga trống trải, hai tay đút trong túi áo khoác, mắt nhìn lên trên trần mái vòm đã nứt nẻ do thời gian, để lại một khoảng không tĩnh mịch đến rợn người.

Không gian quanh tôi im lìm đến mức, nếu không có tiếng gió len qua khe mái tôn rỉ sét, tôi đã tưởng mình vừa điếc. Trời bên ngoài mù xám, tầng mây dày đặc che khuất mặt trời, khiến ánh sáng trong ga chỉ còn lại thứ màu lam tro nhạt nhòa. Cái lạnh luồn sâu qua lớp không khí ẩm, cắn vào da thịt, tôi rụt vai lại, tự cảm ơn rằng ít nhất mình đang diện chiếc áo len và áo khoác từ bộ đồng phục mùa đông của trường Aogami thời hiện đại.

Phía trên cao, những dầm sắt đen sạm đong đưa mỗi khi gió thổi mạnh. Có lúc, tôi nghe tiếng kim loại rên khẽ, như thể chính tòa nhà cũng đang chịu đựng với thời gian cùng tôi.

Tôi ngồi xuống ghế gỗ đã mục nát, hơi nghiêng đầu nhìn tấm biển gỗ cũ kỹ treo lơ lửng: ``Nhà ga Aogami.'' Mấy con chữ khắc tay đã bị thời gian bào mòn, nhưng vẫn còn đủ rõ để nhận ra. Thật kỳ lạ, dù nơi đây hoang tàn, lạnh lẽo, tôi vẫn cảm thấy có gì đó thân quen.

``\eng{Lo sao?}'' giọng Game vang lên bên tai, đều đều như một tần số nhỏ trong đầu tôi.

Tôi khẽ cười, không quay lại: ``\eng{Chắc ông cũng đoán được mà. Nếu nói không lo thì chẳng khác nào tự lừa dối với bản thân cả.}''

``\eng{Thật chứ?}''

``\eng{Tất nhiên rồi.}'' tôi đáp. ``\eng{Không đến mức sợ hãi đâu, nhưng không hẳn là không bồn chồn.}''

Một khoảng im lặng. Tiếng gió rít qua mái tôn thay cho câu trả lời.

``\eng{Hai đứa ấy trông thế thôi mà cũng mạnh mẽ hơn cậu nghĩ đấy,}'' Game nói, giọng nhẹ như hơi thở. ``\eng{Đừng quên, cả hai đã cùng cậu sống sót trong The Void. Chỉ là giờ mỗi người đang phải vượt qua phần của mình thôi.}''

Tôi hít sâu. ``\eng{Biết chứ. Kaigou-chan không dễ gục đâu, cái kiểu sấn sa sấn sổ như cô ấy tôi biết rõ. Còn Suzuki thì\dots}'' tôi dừng lại, ngước nhìn bầu trời xám, thở dài một hơi. ``\eng{Con bé đã trưởng thành hơn tôi tưởng. Từ khi là một cô bé nhút nhát chỉ biết bám lấy hai anh chị, mà giờ đã có thể tự mình giải quyết được rồi. Con bé cũng giúp ích chúng tôi nhiều lắm chứ.}''

Tôi mỉm cười, nhưng nụ cười không hẳn là vui.

``\eng{Thật lạ,}'' tôi nói nhỏ, ``\eng{Ngày xưa, tôi quen việc đi một mình. Không cần ai cả. Nhưng bây giờ\dots\ lại thấy khó chịu khi không có họ ở đây.}''

Game không trả lời. Có lẽ ông hiểu, hoặc cũng chẳng cần nói gì thêm.

Tôi cứ thế ngồi chờ. Thời gian kéo dài như sợi dây cao su, mảnh nhưng căng ra đến đau đớn.

*CỘP CỘP CỘP*

Một tiếng rất khẽ từ phía cánh Tây. Nghe như tiếng bước chân. Không vội vã, không nặng nề, mà như người đang cố giữ lại từng bước đi dưới gót giày để không đánh mất chính mình.

Tôi đứng bật dậy, ánh mắt hướng về phía âm thanh. Từ trong làn sương loãng, bóng dáng một người đang tiến đến, tập tễnh nhưng vững vàng. Một cô gái với mái tóc đuôi ngựa, trông khá xù, kèm với đó là buộc tóc đỏ, và đôi mắt xanh, nhưng đôi mắt ấy không còn là ánh sáng lạnh lùng của Saikyou, mà là biển cả đang rút vào bờ sau cơn bão.

Có lẽ nào đó là\dots

``\textbf{Kaigou-chan!}'' tôi gọi lớn, tiếng gọi vỡ ra thành từng mảnh sương.

Cô ngẩng lên, ánh mắt quen thuộc ánh lên chút mệt mỏi xen lẫn nhẹ nhõm, như một thứ cảm giác lâng lâng sao bao đêm không ngủ. ``\vie{Cậu thực sự\dots\ ở đây à.}''

Tôi gật đầu, tiến lại gần. Mỗi bước chân tôi đi đều cố gắng không gây ra tiếng động, sợ rằng nếu tiếng bước chân quá lớn, cô ấy sẽ tan biến như bong bóng xà phòng. ``\vie{Cô nghĩ tôi sẽ bỏ đi trước chắc?}''

Một nụ cười nho nhỏ thoáng qua môi cô ấy. Hai chúng tôi không cần nói thêm gì, chỉ đơn giản là đứng cạnh nhau, im lặng mà đủ hiểu, kiểu im lặng chỉ những người từng sống sót cùng nhau mới có thể chia sẻ được.

Nhưng chưa kịp thả lỏng, tôi nghe một âm thanh khác, lần này từ phía cánh Đông. Tiếng chạy gấp, tiếng giày va vào nền đá, xen lẫn hơi thở dồn dập, một thứ âm thanh của đứa trẻ đang chạy trốn khỏi cơn ác mộng và chạy thẳng vào vòng tay người thân.

Tôi quay người lại, một bóng dáng nhỏ nhắn đang lao tới, mái tóc đen cột hai bím tung lên theo gió, cùng đôi mắt xanh-đỏ sáng bừng trong màn tối xẩm. Đôi mắt ấy đã từng nhìn thấy cả thế giới sụp đổ qua khe cửa phòng kín, và giờ đang tìm kiếm hai người duy nhất còn sót lại.

``\textbf{Onii-sama! Onee-sama!}'' giọng Suzuki vang lên giữa không khí lạnh buốt, hai tiếng gọi ấy va vào nhau tạo thành một thứ âm thanh nghe như tiếng khóc của con nai con lạc đàn.

Kaigou-chan tròn mắt. ``Suzu\dots?''

Cô bé suýt ngã vì vấp phải một cục gạch, nhưng vẫn lao đến ôm chầm lấy cô chị, rồi ngẩng lên nhìn tôi, nước mắt lăn dài nhưng ánh mắt sáng rực, như thể đó là ánh mắt của đứa trẻ vừa tìm thấy nhà sau bao năm lang thang trong đêm.

``\vie{Em\dots\ em tưởng sẽ không bao giờ gặp lại hai người nữa chứ\dots}''

Tôi khẽ đặt tay lên đầu cô bé, cười nhẹ: ``\vie{Còn anh thì biết chắc rằng em sẽ đến mà.}''

Không ai nói gì thêm trong giây lát. Cơn gió lướt qua, cuốn theo bụi mù và những chiếc lá khô.

Ba chúng tôi, ba người trong cùng họ Hikawa, cuối cùng cũng lại đứng cạnh nhau như ba mảnh vỡ của cùng một tấm gương cuối cùng đã được ghép lại.

Sau khi cả ba đã bình tĩnh, tôi lên tiếng:
``\vie{Có lẽ đã đến lúc chúng ta nói rõ với nhau rồi. Về việc\dots\ ta thực sự là ai.}''

Suzuki gật đầu, mắt vẫn còn sưng đỏ. ``\vie{Em\dots\ nhớ lại hết rồi. Em không còn là Reiko. Em là Hikawa Suzuki. Em là em gái của Onee-sama, đồng thời coi Hikawa Kuon là Onii-sama!}''

Kaigou siết nhẹ vai em. ``\vie{Còn chị\dots}'' cô hít sâu. ``\vie{Chị là Hikawa Kaigou. Không phải Saikyou. Không phải ai khác.}''

Tôi nhìn hai người họ, lòng dâng lên cảm giác rất khó tả.

``\vie{Tốt rồi\dots}'' tôi nói khẽ. ``\vie{Vì anh cũng nhớ lại hết rồi. Anh không còn là Kyuen nữa.}''

Chúng tôi im lặng trong vài giây, như thể đang lắng nghe hơi thở của thế giới quanh mình. Lúc này đây, ánh sáng mờ từ khe mây chiếu xuống, phủ lên ba chúng tôi một lớp vàng nhạt yếu ớt, nhưng ấm áp.

``\vie{Chúng ta\dots}'' tôi nói, nở một nụ cười thật lòng, ``\vie{đã trở lại là chính mình rồi.}''

\begin{center}
  \uline{Ga Aogami? - Sân ga.}
  \pov{???}
\end{center}

\dots

\dots

``\dots Ể?''

\img{e2-8j}

Tôi nhảy xuống con tàu điện đó, khiến tôi bước hụt, suýt chút nữa thì ngã xuống.

Tôi thở hồng hộc, ngực thì loạn xạ hết cả lên.

``Vừa rồi\dots\ là gì vậy? Và\dots\ ể?''

Không còn tiếng bánh sắt nghiến trên đường ray, không còn tiếng rung lắc quen thuộc của con tàu. Trước mắt tôi chỉ là một khoảng trống mênh mông, tĩnh lặng đến rợn người. Tôi\dots\ không còn ở trên tàu nữa.

Thay vào đó, là một nhà ga hoang tàn. Vòm trần vỡ nát, dây điện đứt rủ xuống như những sợi rễ khô. Bụi phủ dày đến mức mỗi bước chân tôi dẫm xuống đều phát ra tiếng kẽo kẹt khô khốc. Không khí đặc quánh, mùi sắt gỉ và tro cũ xen lẫn với hơi lạnh thấm vào da thịt.

Trông nó\dots\ không khác nào là một nhà ga được dựng từ thế kỷ trước, pha trộn với thời đại ngày nay.

Ở phía đối diện, giữa đống gạch vụn và cột bê tông đổ nát, có một tấm bảng gỗ mục nát nghiêng nghiêng, dòng chữ đen sờn đi theo năm tháng: ``Thị trấn Aogami.''

Tôi đứng sững lại, mồ hôi chảy trên má. ``Đ-Đây là\dots''

Không còn nghi ngờ gì nữa, đây chính là thị trấn Aogami, nhưng\dots\ lại ở một thời điểm không xác định.

Không thể nào. Tôi chắc chắn mình vừa rời khỏi con tàu cùng với cô ấy - ``Shino,'' cô gái mang đôi mắt xanh. Tôi vẫn còn nhớ rõ bàn tay ấy đã nắm lấy tôi, kéo tôi đi trong luồng sáng lóa mắt. Thế mà giờ đây, cô ấy đã biến mất.

``Cô ấy đâu rồi\dots?'' tôi lẩm bẩm, bước nhanh về phía sân ga chính. Không một bóng người. Chỉ có tiếng gió rít qua những khe gạch vỡ.

Một tia sáng phản chiếu nơi góc tường khiến tôi dừng lại. Là một tấm gương nứt toác, có lẽ từng thuộc về phòng chờ hành khách. Tôi cúi xuống nhìn, và toàn thân đông cứng lại.

Trong gương\dots\ là cô ấy.

Mái tóc dài màu nâu, chiếc cài tóc lệch ở trên và làn da trắng bệch, đi cùng với bộ đồng phục màu xanh đậm và khăn quàng màu xanh dương kia\dots\ không thể nhầm vào đâu được. Màu mắt của tôi lúc này đây\dots\ là màu xanh nước biển, không còn màu nâu nữa.

\img{e2-8k}

Tôi há hốc miệng, cổ họng nghẹn cứng. ``Kh-Không thể nào\dots\ Tại sao\dots\ cô ấy lại ở trong hình phản chiếu của mình chứ\dots?''

Rồi trong khoảnh khắc ấy, tim tôi như ngừng đập.
Không. Không phải ``cô ấy'' ở trong gương.
Là tôi.

``T-Tại sao lại là mình\dots!? Không lẽ\dots\ mình đang ở trong cơ thể của cô ấy!?''

Tôi giật lùi, hơi thở dồn dập. Những câu hỏi không ngừng gào thét trong đầu tôi: ``Tại sao?'' ``Khi nào?'' ``Làm thế nào!?''

Tôi quay người, lao về phía trung tâm sân ga, nơi mà tôi mơ hồ cảm nhận được một sự hiện diện khác. Nhưng đôi chân tôi run rẩy, chẳng còn cảm giác đâu là thật nữa.

``Tại sao\dots\ \textbf{tại sao chứ\dots!?}'' tôi vừa chạy vừa gào lên, nước mắt nóng rát nơi khóe mắt.

Mỗi nhịp tim vang trong đầu như tiếng đập cửa, thúc giục tôi phải tìm ra câu trả lời. Nhưng càng chạy, khung cảnh xung quanh càng méo mó, mọi thứ cứ xoay tròn, cuốn tôi vào một cơn lốc mù mịt.

Và tôi nhận ra: tôi không còn là chính mình nữa.

Không biết mình là Sora\dots\ hay Shino.

Không biết đâu là thật, đâu là giấc mộng.

Chỉ biết rằng, phần ``tôi'' đang dần tan biến trong cái xác xa lạ này, như thể cả hai linh hồn đang hòa lẫn vào nhau, đến mức chẳng còn ranh giới nào để quay lại.

Tôi thở hổn hển, bàn tay run rẩy chạm vào ngực mình. Tim vẫn đập, mạnh và đều, nhưng không phải của tôi. Nó đập bằng nhịp điệu của \emph{cô ấy}.

Không\dots\ không thể nào. Những mảng âm thanh méo mó còn sót lại từ toa tàu phía sau vẫn còn văng vẳng ở đâu đó. Mọi thứ quanh tôi như đang xoay tròn, rồi dần ổn định lại thành hình ảnh của một nhà ga đổ nát.

Bước chân tôi dừng lại đột ngột khi nghe thấy tiếng người vọng lại phía bên trong nhà ga.

``\vie{Suzuki, em không xây xát ở chỗ nào đấy chứ?}'' giọng một người con trai cất lên.

``\vie{Em ổn ạ\dots\ Onii-sama,}'' giọng một cô gái cất lên, có chút gì đó hơi cao vút.

``\vie{Cô thì sao, Kaigou-chan?}''

``\vie{Không việc gì. Tôi nghĩ tôi phải hỏi cậu câu ấy đó, Kuon-kun,}'' một giọng nữ khác cất lên, nhưng có gì đó trầm tĩnh và\dots\ mạnh mẽ hơn hẳn.

Tôi thoáng giật mình. Kuon? Kaigou? Suzuki?

Ba giọng nói này\dots\ nghe vừa quen thuộc, lại vừa cảm thấy lạ.

Những cái tên đó đâm thẳng vào đầu tôi như một tia sét.

Không phải Kyuen, Saikyou hay Reiko như tôi đã từng gọi họ.

Không phải cái tên của những người mà tôi đã quen biết lúc đó.

Tôi lùi lại nửa bước, tim đập loạn.

``Không thể nào\dots\ Tên của họ\dots\ sao lại\dots''

Họ ngẩng lên gần như cùng lúc. Ba ánh mắt đổ dồn về phía tôi. Kuon hơi nhíu mày, Kaigou siết chặt tay, còn Suzuki chớp mắt đầy hoài nghi.

``\dots Shino-san?'' cậu con trai, người mà tôi từng gọi là ``Kyuen'' cất tiếng.

``Kuon-kun, tôi không nghĩ đó là Shino-san đâu,'' cô gái có đôi mắt xanh nói. Cả cách xưng hô của cô với cậu ấy cũng khác hẳn.

``Sao em thấy\dots\ có gì đó quen thuộc ghê,'' cô bé có mắt hai màu tiếp lời. ``Như thể là cậu ấy với Sora-san đã nhập vào một vậy.''

Tôi sững sờ. Quả thật\dots\ Reiko-san - hoặc ít nhất, đã từng như vậy - đoán đúng. Thật sự tôi vẫn không thể hiểu nổi chuyện gì đang xảy ra với tôi nữa.

``Cậu là\dots\ Kuon? Còn cậu là\dots\ Suzuki? Còn cô\dots'' tôi nhìn cô gái cao lớn kia, ``\dots Kaigou\dots?''

Họ không đáp, chỉ gật nhẹ, ánh mắt vẫn giữ nguyên vẻ cảnh giác.

``Nhưng\dots\ tại sao?'' tôi nói, giọng run như sắp vỡ. ``Tại sao mọi người\dots\ lại mang những cái tên đó? Tôi nhớ rõ\dots\ đó không phải là các cậu mà mình thấy\dots''

Tôi ôm đầu, lùi lại, hơi thở dồn dập. ``Không thể nào\dots\ Không thể nào\dots''

Và rồi, loa phát thanh trên cao lại vang lên, tiếng rè đặc trưng xen lẫn tiếng thở khàn của một giọng nam đậm âm sắc Pháp: ``Ồ, vậy ra đó là Shino-san mà Kuon-user đã bảo đó sao.''

Tôi quay ngoắt đầu, mắt đảo quanh. ``Ông là ai!?''

``Cứ gọi tôi là Game,'' giọng nói đáp, nghe như đang cười. ``Cô có vẻ hơi rối nhỉ.''

``Rối!? Tất nhiên là rối rắm rồi! \textbf{Cơ thể này không phải của tôi! Tôi không biết mình là ai nữa!}'' tôi hét, giọng khản đặc.

Một tràng rè nhẹ vang lên trước khi ông ta lại nói, lần này nhẹ giọng hơn: ``Bình tĩnh. Cảm giác `không phải mình' ấy\dots\ là do hai phần ký ức đang chồng chéo nhau thôi. Một phần là của cô, phần kia là của Shino thật.''

``Shino thật\dots?''

``Cô là Sora, phải không? Còn cơ thể này, là của Misora Shino, cô gái đã từng tồn tại ở thị trấn này. Hai linh hồn cùng ở trong một vỏ xác\dots\ Tất nhiên là sẽ khó chịu rồi.''

Tôi lắc đầu quầy quậy, nước mắt trào ra lúc nào không hay. ``Vậy nghĩa là tôi đã cướp lấy cô ấy\dots?''

``Không. Cô không cướp gì cả. Cô chỉ nhập lại thôi.''

``Nhập lại\dots?''

``Hai người vốn là một ngay từ đầu. Hai mảnh của cùng một linh hồn. Một người đã vốn đã chấp nhận với sự thật, một người đã ruồng bỏ trốn chạy khỏi nó. Giờ chỉ là\dots\ cả hai đang gặp lại nhau thôi.''

Tôi im lặng. Những lời đó lướt qua tâm trí tôi, để lại dư âm kỳ lạ: vừa đau đớn tột cùng, nhưng lại vừa cảm thấy nhẹ nhõm hơn hẳn.

Ở phía xa, tôi thấy ba người họ vẫn đứng đó, ánh mắt họ lộ rõ vẻ bối rối, nhưng không ai tiến lại. Có lẽ, họ cũng đã nhận ra có điều gì không ổn ở tôi.

Tôi khẽ nói, giọng run rẩy: ``\dots Tôi không biết mình là ai nữa. Nhưng\dots\ tôi biết chắc rằng, những người trước mặt tôi\dots\ họ là thật. Không phải ảo ảnh, không phải ký ức.''

Cậu con trai bước lên một bước, ánh mắt sắc lạnh nhưng không mất đi sự cảm thông. ``Có lẽ\dots\ việc Sora-san đi theo Shino-san chính là cách mà hai mảnh linh hồn ấy hợp nhất. Một người đã sẵn sàng đối mặt, một người vẫn còn do dự. Nhưng khi họ cùng bước qua cánh cửa kia\dots\ họ không còn là hai người nữa.''

Tôi lặng đi. Hai mảnh linh hồn? Tôi chưa từng nghĩ đến việc Sora và Shino có thể là một. Nhưng giờ nghe cậu ấy nói, tôi chợt thấy từng mảnh ghép trong đầu mình khớp lại: ánh mắt họ khi nhìn nhau, cách họ né tránh, cách họ tìm kiếm. Tôi đã từng nghĩ họ là hai cực đối lập, nhưng có khi họ chỉ là hai mặt của cùng một đồng xu đã bị bẻ gãy.

Kaigou-san gật đầu, giọng khẽ nhưng chắc nịch: ``Tôi đã thấy ánh mắt cậu trước khi rời đi. Không còn là Sora của những ngày đầu nữa. Cậu không còn né tránh. Cậu ấy\dots\ đã chọn hướng về sự thật.''

Tôi thấy lòng mình nhói lên. Chọn ư? Chọn gì? Chọn buông tay? Chọn đối mặt? Hay chọn trở thành một phần của Shino? Tôi không biết, nhưng tôi hiểu rằng Sora đã bước qua một ranh giới mà tôi chưa từng dám đến gần. Và tôi, kẻ đứng đây, vẫn còn mắc kẹt giữa những mảnh vỡ của chính mình.

Suzuki-san nắm chặt tay áo, giọng nhỏ nhẹ nhưng chắc chắn, củng cố với luận điểm của hai người. ``Và Shino-san cũng vậy. Cậu không còn đứng một mình. Cô ấy đã đợi điều này từ lâu lắm rồi. Để được gặp lại\dots\ phần còn lại của mình.''

Tôi nuốt khan. Phần còn lại của mình. Tôi chợt tự hỏi: còn tôi thì sao? Tôi có còn phần nào đang chờ được gặp lại không? Hay tôi đã vứt bỏ nó từ lâu, trong một vòng lặp nào đó, khi tôi quyết định trở thành kẻ chỉ biết chạy trốn? Tôi không có câu trả lời. Chỉ biết rằng, khi họ nói vậy, tôi thấy mình nhỏ bé đến lạ thường.

Loa rè lên lần cuối, Game cười nhỏ: ``Vậy thì tốt rồi. Bởi vì, chỉ những ai vẫn còn tin vào `thật' mới có thể bước tiếp trong thế giới này.''

Tôi siết chặt tay. Tin vào ``thật'' ư? Tôi đã từng tin. Nhưng giờ đây, sau bao nhiêu lần tháo chạy, tôi không còn chắc điều gì là ``thật'' là gì nữa. Có phải là những gì tôi đang thấy? Hay là những gì tôi đã cố quên? Tôi không biết. Nhưng tôi biết một điều: nếu tôi không bước tiếp, tôi sẽ mãi mất họ. Và có lẽ, mất luôn cả chính mình.

Tín hiệu tắt ngúm. Tiếng rè biến mất, để lại trong nhà ga khoảng lặng nặng nề.

Tôi hít sâu, nhìn ba người trước mặt. Lòng bất an cuộn lên như cơn sóng ngầm. Tay tôi siết lại đến mức đầu móng tái dại.

Câu nói của ông ta chẳng khác nào như một bản án, và tôi biết bản thân mình đang trả giá bằng chính ký ức đã bị cắt xén.

Tôi đứng đó, nhìn ba người họ bước về phía ánh sáng trắng, mỗi bước chân như giẫm lên từng mảnh vỡ trong lòng tôi. Nhưng rồi, một tiếng nói vang lên, không phải từ ai xa lạ, mà từ chính Kuon, người mà tôi từng nghĩ rằng đã bị tôi đẩy vào quên lãng.

``Shino-san,'' cậu ta gọi tên tôi, giọng không giận dữ, không trách móc, nhưng cũng có một cái gì đó nghiêm nghị. ``Giống như bản ngã kia của cậu từng nói đấy\dots\ cậu vẫn còn muốn chạy đến bao giờ?''

Tôi cứng đờ. Câu hỏi đó quá quen thuộc, quen đến mức tôi đã từng tự hỏi chính mình điều đó hàng trăm lần trong những vòng lặp đã qua. Nhưng tôi chưa từng dám trả lời.

Kaigou bước lên, đứng cạnh cậu ấy. ``Chúng tôi không trở lại để trách cậu. Chúng tôi trở lại vì\dots\ chúng tôi vẫn còn tin rằng cậu có thể dừng lại.''

Tôi mở miệng, nhưng không thốt nên lời. Cổ họng tôi như bị bóp nghẹt bởi hàng triệu mảnh ký ức mà tôi từng cố gắng xóa đi, những khuôn mặt tôi đã lừa dối, những lời hứa tôi đã không giữ, những vòng lặp tôi đã tạo ra chỉ để níu kéo một thế giới không còn thật.

Suzuki bước đến gần nhất. Cô đưa tay ra, đơn giản chỉ\dots\ chạm vào vai tôi, nhẹ như một lời thì thầm.

``Cậu đã cố gắng rồi,'' cô ấy nói, giọng nhỏ mà rõ ràng. ``Nhưng giờ thì\dots\  đủ rồi.''

Tôi run rẩy. Không phải vì sợ, mà vì\dots\ tôi đang khóc. Lần đầu tiên trong hàng trăm vòng lặp, tôi không khóc vì tuyệt vọng, mà vì tôi được nhìn thấy họ lần nữa. Không như những con rối trong thế giới tôi dựng lên, họ là những con người thật sự, vẫn còn đó, vẫn còn nhìn thấy tôi, và vẫn chừa cho tôi một con đường lui.

``Mình\dots\ mình đã nghĩ rằng nếu mình không dừng lại, thì mọi thứ sẽ không kết thúc,'' tôi nói, giọng nghẹn ngào. ``Nhưng mình đã lầm. Mình chỉ\dots\ kéo dài nỗi đau của chính mình, và của cả những người\dots\ mà mình đã tôn trọng và yêu thương.''

Tôi cắn chặt môi, mặt nóng bừng. ``Mình đã quá ích kỷ\dots\ quá sa đà vào mong ước của riêng mình, đến mức quên rằng họ cũng đau, cũng mệt, cũng muốn được nghỉ. Mình biến họ thành con rối, biến ký ức thành xiềng xích, chỉ để níu kéo một giấc mơ đã thối rữa. Và rồi\dots\ mọi thứ vượt khỏi tầm tay, như con tàu lao dốc không phanh, đâm thẳng vào bóng tối.''

Kuon gật đầu, như thấu hiểu được những gì mà tôi đang nghĩ. ``Chúng tôi cũng biết cậu đã phải chịu khổ quá nhiều. Nhưng chúng tôi không muốn cậu tiếp tục mang nỗi đau ấy một mình.''

Tôi nhìn họ,  ba người mà tôi từng nghĩ rằng đã mất vĩnh viễn. và lần đầu tiên, tôi không thấy họ như những phần thưởng cho sự kiên trì của tôi, mà như những người đang chờ tôi buông tay.

Tôi biết, rằng khi bước qua ánh sáng trắng kia, tôi sẽ không còn có thể quay lại. Tôi sẽ phải đối mặt với tất cả những gì tôi đã làm, không chỉ trong thế giới này, mà trong chính trái tim tôi.

Nhưng lần này\dots\ tôi không còn đi một mình.

Tôi biết rất rõ, rằng ở bên kia ánh sáng trắng ấy sẽ là một sự thật tàn nhẫn mà tôi sẽ phải đón nhận.

Mỗi bước họ tiến về phía ánh sáng, tôi lại cảm thấy một mảnh ghép trong lòng mình bị bóc ra, rơi xuống vực thẳm của vòng lặp.

Bởi vì\dots\ chính tôi là người đã từng là là kẻ đẩy họ vào mớ rắc rối này mà.

\pov{-}

Gió thổi mạnh hơn khi cánh cửa ga tàu khép lại sau lưng họ. Ánh sáng ngoài kia phủ xuống như một tấm màn xám bạc, bao trùm lên cảnh tượng đổ nát đến rợn người của thị trấn Aogami. Con đường vốn từng đông đúc giờ chỉ còn lại những tòa nhà nứt vỡ, dây điện đổ rạp, biển hiệu gãy gập treo lủng lẳng trong gió. Không còn tiếng người, không còn tiếng động cơ, chỉ còn tiếng gió rít qua những khung cửa vỡ và tiếng sắt thép nghiến vào nhau như tiếng than khóc mệt mỏi của một vùng đất đã chết.

Kuon dừng lại trước ngã tư. Dưới chân anh, từng tờ báo cũ bị gió cuốn bay loạn xạ. Một vài tờ vướng vào hàng rào gãy, vài tờ khác bị dính trong vũng nước mưa đã chuyển màu nâu đục. Kaigou cúi xuống nhặt một mảnh còn sót lại, nhưng nó rách nát đến mức chỉ còn vài dòng chữ mờ.

``Ngày X tháng 7 năm 1946 - Thảm kịch tại Trường Trung học Aogami\dots''

Cậu chưa kịp đọc hết thì gió lại cuốn tờ giấy đi, biến nó thành một đốm trắng lẫn vào bụi mù xa xa.

Shino bước chậm phía sau. Khi cô nhìn xuống chân, một mẩu báo khác nằm đó, nhàu nát, lem nước, nhưng rõ ràng vẫn còn thấy được hình ảnh đen trắng của vài gương mặt học sinh. Cô dừng lại, lặng nhìn. Những gương mặt ấy dường như đang mờ dần đi, như thể bị chính thời gian và ký ức nuốt chửng.

Suzuki đến gần, khẽ nói: ``Có lẽ\dots\ đó là những người đã mất trong vụ sạt lở năm ấy. Trong đó\dots\ có cả Shino-san nữa.''

Không ai trả lời. Chỉ có gió tiếp tục thổi, mang theo những tờ báo bay lên không trung như đàn bướm tàn tro.

Bốn người họ cùng ngẩng đầu nhìn về phía xa, nơi những đám mây đen phủ kín bầu trời, và phía chân trời, ánh sáng của hoàng hôn trộn lẫn cùng tro bụi tạo nên một sắc đỏ kỳ lạ. Trong khoảnh khắc ấy, họ đều cùng hiểu một điều, rằng mọi thứ họ đang thấy chỉ là phần nổi của một điều gì đó sâu hơn, tối tăm hơn đang ẩn dưới lớp đất mục nát này.

Kaigou siết chặt nắm tay, Kaigou giương ánh mắt quyết tâm về phía trước Shino gật đầu còn Suzuki tỏ vẻ một chút lo lắng.

``Có lẽ\dots'' Kuon khẽ nói,  ``chúng ta đã tiến gần hơn rồi.''

Cô nàng ``song sinh'' mở mắt, nghiêng đầu: ``Ý cậu là\dots?''

``Tới sự thật.''

Cậu đáp, ánh mắt hướng về phía tấm biển gỗ ``Aogami'' bên ngoài nhà ga đã xói mòn vì thời gian.  ``Mỗi toa tàu, mỗi giấc mơ, mỗi thứ chúng ta thấy\dots\ tất cả đều đang dẫn đến cùng một điểm. Có thể, thị trấn này chưa bao giờ là thật. Nhưng cũng có thể, nó là thật theo một cách nào đó mà ta chưa hiểu được.''

Cô em út ngẩng lên, giọng nhẹ như gió: ``Vậy chúng ta sẽ làm gì bây giờ ạ?''

Kuon nhìn sang hai người, khẽ mỉm cười. ``Chúng ta vẫn làm điều mà mình luôn làm.''

Cô bé chớp mắt. ``Là gì ạ?''

``Tiếp tục cùng nhau hướng về phía trước.''

Kaigou-chan bật cười khẽ, tiếng cười hiếm hoi sau bao nhiêu hỗn loạn. ``Cậu nói thì dễ như bỡn lắm nhỉ? Tổng quản có khác!''

``Vì tôi tin vậy mà,'' Kuon đáp, rồi nhìn thẳng vào cô ấy.

``Miễn là chúng ta có nhau, chúng ta sẽ tìm được cách để giải mã bí ẩn này. Dù thị trấn Aogami này có là mảnh ký ức, hay một vũ trụ méo mó nào đó, thì vẫn sẽ có chìa khóa. Và tôi tin\dots''

Kuon dừng lại, nắm chặt tay lại. ``\dots chìa khóa ấy nằm trong tay chúng ta.''

Shino bước lên, đứng cạnh Kaigou, khẽ nói: ``Nếu cậu đã tin vậy, thì mình cũng sẽ tin theo. Dù chìa khóa có nằm ở đâu, mình sẽ cùng cậu tìm ra.''

Kaigou gật đầu, nắm lấy tay Shino: ``Chúng ta cùng nhau, được chứ?''

Suzuki mím môi, rồi cũng gật đầu, giọng nhỏ nhưng rõ ràng: ``Được ạ. Dù có chuyện gì xảy ra, em cũng sẽ đi cùng mọi người.''

Gió lại thổi qua nhà ga, mang theo âm thanh mơ hồ của tiếng chuông tàu vọng lại từ xa, như một lời nhắc nhở rằng hành trình vẫn chưa kết thúc.

Bốn người họ nhìn nhau, không ai nói gì thêm. Chỉ có một cảm giác chung, rằng dù con đường phía trước vẫn bị bao phủ trong sương mù, thì ít nhất, lần đầu tiên sau bao nhiêu ngày lạc lối, họ đã biết mình đang đi đúng hướng.
Một bước nữa thôi, sự thật của thị trấn Aogami sẽ lộ diện.

Dù thị trấn này có bị chôn vùi trong quá khứ, dù những bóng ma vẫn lẩn khuất đâu đó giữa tàn tích, họ vẫn sẽ tiếp tục bước đi.

Miễn là họ còn ở bên nhau, họ sẽ tìm được cách để giải mã bí ẩn này.

Và đâu đó, trong cơn gió lạnh lẽo của hoàng hôn, vang lên tiếng chuông tàu điện xa xăm, như lời nhắc rằng hành trình của họ vẫn chưa kết thúc.

Bởi có lẽ, chính sự thật của Aogami mới là chìa khóa để giải mã những bí ẩn xung quanh nơi này.

\end{document}
